\chapter[Affine Toric Varieties]
{Affine Toric Varieties:\\ Semigroups, Cones, and Limits}
\label{chap:affine-toric-varieties}

\begin{chapteroverview}
Which monomial functions on a torus remain regular after we add a
boundary?

For the affine plane, the answer is familiar: the Laurent monomial
\(x^ay^b\) extends precisely when \(a,b\geq0\). This chapter turns that
example into a classification. Additive collections of characters
produce semigroup algebras and affine varieties; normality turns those
collections into the lattice points of rational polyhedral cones.

The same variety then admits two dual descriptions:
\[
\begin{array}{c}
\text{a cone in \(M_{\RR}\) records the characters that extend,}
\\[1mm]
\text{while a cone in \(N_{\RR}\) records the one-parameter
subgroups whose limits exist.}
\end{array}
\]
Our main goal is to understand both readings of
\[
U_\sigma
=
\Spec\Bbbk[\sigma^\vee\cap M].
\]

When a result has a lengthy or technical proof, we will give a precise
reference to Cox--Little--Schenck and a conceptual proof sketch that
isolates the driving idea and the logical flow.
\end{chapteroverview}

\begin{threedictionaries}
Affine toric geometry begins with one object viewed in three ways.

\medskip

\begin{center}
\small
\begin{tabularx}{\linewidth}{@{}
  >{\raggedright\arraybackslash}X
  >{\raggedright\arraybackslash}X
  >{\raggedright\arraybackslash}X@{}}
\toprule
\textbf{Algebra}
&
\textbf{Geometry}
&
\textbf{Combinatorics}
\\
\midrule
a semigroup algebra \(\Bbbk[S]\)
&
an affine toric variety \(\Spec\Bbbk[S]\)
&
an additive semigroup \(S\subseteq M\)
\\[1mm]
regular characters
&
monomial functions extending to the boundary
&
lattice points in a cone in \(M_{\RR}\)
\\[1mm]
pullback along \(\oneps{u}\)
&
limits of one-parameter motions
&
a dual cone in \(N_{\RR}\)
\\
\bottomrule
\end{tabularx}
\end{center}
\end{threedictionaries}

\section{From the affine plane to affine toric varieties}
\label{sec:affine-plane-to-affine-toric-varieties}
\index{affine toric variety}
\index{dense torus}
\index{character!extending to an affine variety}

We begin where Chapter~0 began:
\[
\AAA^2=\Spec\Bbbk[x,y].
\]
The affine plane contains the torus
\[
T=\Gm^2
=
\{(x,y)\in\AAA^2:xy\neq0\}
\]
as a dense open subset. On \(T\), both \(x\) and \(y\) are invertible,
so every Laurent monomial
\[
x^ay^b,
\qquad
(a,b)\in\ZZ^2,
\]
is a regular function. In the language of Chapter~1, these Laurent
monomials are precisely the characters of \(T\):
\[
x^ay^b=\character{(a,b)}.
\]

But most of these characters do not remain regular on all of
\(\AAA^2\). The function \(x^{-1}\), for instance, has a pole along
the coordinate line \(V(x)\). Likewise, \(y^{-1}\) has a pole along
\(V(y)\). The characters that do extend regularly are exactly
\[
x^ay^b,
\qquad
a,b\geq0.
\]
Thus
\[
\Bbbk[\AAA^2]
=
\Bbbk[x,y]
=
\bigoplus_{(a,b)\in\ZZ_{\geq0}^2}
\Bbbk\,\character{(a,b)}.
\]

This gives us the question that drives the entire chapter:

\begin{guidingquestion}
Given a torus \(T\), which collections of characters can occur as the
regular functions on an affine variety containing \(T\)?
\end{guidingquestion}

\subsection{Which characters extend?}

For the affine plane, the answer is visible in the character lattice
\[
M\simeq\ZZ^2.
\]
The lattice point \((a,b)\) corresponds to the character \(x^ay^b\),
and this character extends to \(\AAA^2\) exactly when
\[
(a,b)\in\ZZ_{\geq0}^2.
\]
The allowed exponent vectors form the lattice points of the first
quadrant in \(M_{\RR}\).

Two nearby examples help separate the roles of affine and invertible
coordinates.

\[
\begin{array}{c|c|c}
X
&
\Bbbk[X]
&
\text{characters that extend}
\\
\hline
\Gm^2
&
\Bbbk[x^{\pm1},y^{\pm1}]
&
(a,b)\in\ZZ^2
\\[1mm]
\AAA^1\times\Gm
&
\Bbbk[x,y^{\pm1}]
&
a\geq0,\quad b\in\ZZ
\\[1mm]
\AAA^2
&
\Bbbk[x,y]
&
a\geq0,\quad b\geq0.
\end{array}
\]

\begin{figure}[ht]
\centering
\begin{tikzpicture}[scale=0.82]
  % Panel 1
  \begin{scope}[xshift=-6.2cm]
    \fill[DictionaryBlue!8] (-2.1,-2.1) rectangle (2.1,2.1);
    \draw[->,black!55] (-2.25,0)--(2.35,0);
    \draw[->,black!55] (0,-2.25)--(0,2.35);
    \foreach \a in {-2,-1,0,1,2}{
      \foreach \b in {-2,-1,0,1,2}{
        \fill[DictionaryBlue] (\a,\b) circle (1.65pt);
      }
    }
    \node[font=\bfseries] at (0,2.85) {\(\Gm^2\)};
    \node[align=center] at (0,-2.85)
      {all lattice points\\all Laurent monomials};
  \end{scope}

  % Panel 2
  \begin{scope}
    \fill[PastelYellow] (0,-2.1) rectangle (2.1,2.1);
    \draw[->,black!55] (-2.25,0)--(2.35,0);
    \draw[->,black!55] (0,-2.25)--(0,2.35);
    \foreach \a in {0,1,2}{
      \foreach \b in {-2,-1,0,1,2}{
        \fill[ExampleGold] (\a,\b) circle (1.65pt);
      }
    }
    \node[font=\bfseries] at (0,2.85) {\(\AAA^1\times\Gm\)};
    \node[align=center] at (0,-2.85)
      {\(a\geq0\), while \(b\)\\may have either sign};
  \end{scope}

  % Panel 3
  \begin{scope}[xshift=6.2cm]
    \fill[SurfaceGreen!10] (0,0) rectangle (2.1,2.1);
    \draw[->,black!55] (-2.25,0)--(2.35,0);
    \draw[->,black!55] (0,-2.25)--(0,2.35);
    \foreach \a in {0,1,2}{
      \foreach \b in {0,1,2}{
        \fill[SurfaceGreen] (\a,\b) circle (1.65pt);
      }
    }
    \node[font=\bfseries] at (0,2.85) {\(\AAA^2\)};
    \node[align=center] at (0,-2.85)
      {\(a,b\geq0\)\\ordinary monomials};
  \end{scope}
\end{tikzpicture}
\caption{The lattice points indexing regular characters depend on
which boundary directions have been added to the torus.}
\label{fig:characters-extending-basic-affine-examples}
\end{figure}

\begin{pauseandthink}
Consider
\[
X=\AAA^2\times\Gm
\]
with dense torus \(\Gm^3\). Which characters
\[
\character{(a,b,c)}=x^ay^bz^c
\]
extend regularly to \(X\)? Draw the corresponding lattice points in
\(M_{\RR}\simeq\RR^3\).
\end{pauseandthink}

The pattern is already suggestive. Whenever two characters extend,
their product extends as well:
\[
\character{m}\character{m'}
=
\character{m+m'}.
\]
The constant character \(\character{0}=1\) also extends. Therefore the
set of extending characters is closed under addition and contains
\(0\). In the next section, we will call such a collection a
\emph{semigroup}.

\subsection{The torus action on the affine plane}

The dense torus does more than sit inside \(\AAA^2\). It acts on the
whole affine plane by coordinatewise multiplication:
\[
(t_1,t_2)\cdot(x,y)
=
(t_1x,t_2y).
\]
On the open subset where \(xy\neq0\), this is simply the group law of
\(T=\Gm^2\). Thus the multiplication action of the torus on itself
extends across the coordinate axes.

The coordinate axes record what has been added to the torus:
\[
\AAA^2
=
T
\sqcup
\bigl(V(x)\setminus\{(0,0)\}\bigr)
\sqcup
\bigl(V(y)\setminus\{(0,0)\}\bigr)
\sqcup
\{(0,0)\}.
\]
Each piece is a torus orbit. The two one-dimensional orbits are
obtained by allowing one coordinate to vanish, and the origin is the
unique fixed point.

\begin{figure}[ht]
\centering
\begin{tikzpicture}[scale=1.05]
  \draw[->,black!55] (-3.0,0)--(3.15,0) node[right] {\(x\)};
  \draw[->,black!55] (0,-2.7)--(0,2.85) node[above] {\(y\)};

  \draw[very thick,DictionaryBlue] (-2.8,0)--(-0.15,0);
  \draw[very thick,DictionaryBlue] (0.15,0)--(2.8,0);
  \draw[very thick,SurfaceGreen] (0,-2.5)--(0,-0.15);
  \draw[very thick,SurfaceGreen] (0,0.15)--(0,2.5);
  \fill[WarningRed] (0,0) circle (3pt);

  \node[fill=PastelYellow,rounded corners,inner sep=4pt]
    at (1.55,1.35) {\(T=\Gm^2\)};
  \node[DictionaryBlue,below] at (1.65,-0.05)
    {\(V(y)\setminus\{0\}\)};
  \node[SurfaceGreen,left,rotate=90] at (-0.08,1.45)
    {\(V(x)\setminus\{0\}\)};
  \node[WarningRed,below left] at (-0.05,-0.05)
    {\((0,0)\)};
\end{tikzpicture}
\caption{The affine plane consists of its dense torus, two
one-dimensional boundary orbits, and one fixed point.}
\label{fig:a2-torus-orbits-preview}
\end{figure}

The orbit decomposition will eventually be read directly from cones.
For now, it simply illustrates the two features that matter:

\[
\boxed{
\text{\(\AAA^2\) contains \(T\) densely, and the action of \(T\)
extends to all of \(\AAA^2\).}
}
\]

\subsection{The intrinsic definition}

The affine plane suggests the general definition.

\begin{definition}[Affine toric variety]
\label{def:affine-toric-variety}
An \emph{affine toric variety} is an irreducible affine variety \(X\)
containing an algebraic torus \(T\) as a dense open subset such that
the multiplication action
\[
T\times T\longrightarrow T
\]
extends to an algebraic action
\[
T\times X\longrightarrow X.
\]
The torus \(T\) is called the \emph{dense torus} or \emph{big torus}
of \(X\).
\end{definition}

\noindent
\textit{Reference.}
This is Cox--Little--Schenck, Definition~1.1.3, with \(\CC\) replaced
by our algebraically closed base field \(\Bbbk\).

Each condition has a role.

\begin{itemize}
\item \textbf{Affine} means that \(X\) is determined by its coordinate
ring:
\[
X=\Spec\Bbbk[X].
\]

\item \textbf{Irreducible} ensures that the dense torus lies in one
geometric piece and that regular functions on \(X\) are determined by
their restriction to \(T\).

\item \textbf{Dense and open} means that \(X\) is obtained from \(T\)
by adding boundary points, while retaining the torus as the generic
part of the variety.

\item \textbf{The action extends} means that the added boundary is
compatible with the group structure of \(T\). This condition is what
forces monomials and lattice combinatorics to control the geometry.
\end{itemize}

\begin{example}
The torus \(T\) itself is an affine toric variety. Its boundary is
empty, and
\[
\Bbbk[T]=\Bbbk[M].
\]
\end{example}

\begin{example}
Affine space \(\AAA^n\) is an affine toric variety with dense torus
\[
\Gm^n
=
\AAA^n\setminus V(x_1\cdots x_n).
\]
The action is coordinatewise multiplication, and
\[
\Bbbk[\AAA^n]
=
\bigoplus_{m\in\ZZ_{\geq0}^n}
\Bbbk\,\character{m}.
\]
\end{example}

\begin{example}
More generally,
\[
\AAA^r\times\Gm^{n-r}
\]
is an affine toric variety with dense torus \(\Gm^n\). Its coordinate
ring is
\[
\Bbbk[x_1,\ldots,x_r,
x_{r+1}^{\pm1},\ldots,x_n^{\pm1}],
\]
and its extending characters are indexed by
\[
\ZZ_{\geq0}^r\times\ZZ^{n-r}\subseteq M.
\]
\end{example}

\begin{warningbox}
Normality is not part of \cref{def:affine-toric-variety}. In these
notes, as in CLS Chapter~1, affine toric varieties may be nonnormal.
The distinction is essential: later in this chapter, normality will be
recognized combinatorially as a condition that fills in missing
lattice points.

Some authors build normality into the definition of a toric variety.
Whenever conventions could matter, we will state them explicitly.
\end{warningbox}

\subsection{The coordinate-ring viewpoint}

Let \(X\) be an affine toric variety with dense torus \(T\) and
character lattice \(M\). Restricting regular functions from \(X\) to
the dense open subset \(T\) gives an injective homomorphism
\[
\Bbbk[X]\hookrightarrow\Bbbk[T]=\Bbbk[M].
\]
It is injective because a regular function on an irreducible variety
that vanishes on a dense open subset vanishes everywhere.

Thus we may regard \(\Bbbk[X]\) as a subalgebra of the Laurent
polynomial algebra \(\Bbbk[M]\). In concrete terms:

\[
\boxed{
\Bbbk[X]
=
\{\text{regular functions on \(T\) that extend to \(X\)}\}.
}
\]

The torus action places a strong restriction on this subalgebra. The
representation theory of Chapter~1 suggests that it should decompose
into character spaces. Consequently, the coordinate ring should be
determined by the set
\[
S_X
=
\{m\in M:\character{m}\text{ extends regularly to }X\}.
\]
For \(\AAA^2\), this set is \(\ZZ_{\geq0}^2\).

The precise theorem—that every affine toric coordinate ring is a
semigroup algebra \(\Bbbk[S]\)—is Cox--Little--Schenck,
Theorem~1.1.17. We will build and explain that equivalence in the next
section rather than use it as a black box here.

\begin{threedictionaries}
The affine-plane example already contains the basic dictionary.

\medskip

\begin{center}
\small
\begin{tabularx}{\linewidth}{@{}
  >{\raggedright\arraybackslash}X
  >{\raggedright\arraybackslash}X
  >{\raggedright\arraybackslash}X@{}}
\toprule
\textbf{Algebra}
&
\textbf{Geometry}
&
\textbf{Combinatorics}
\\
\midrule
\(\Bbbk[X]\subseteq\Bbbk[M]\)
&
an affine variety \(X\) containing \(T\)
&
a collection of lattice points in \(M\)
\\[1mm]
characters that extend regularly
&
monomial functions without poles on the boundary
&
allowed exponent vectors
\\[1mm]
multiplication of characters
&
multiplication of regular functions
&
addition of lattice points
\\[1mm]
the \(M\)-grading of the coordinate ring
&
the action of the dense torus
&
weights indexed by \(M\)
\\
\bottomrule
\end{tabularx}
\end{center}
\end{threedictionaries}

\begin{guidingquestion}
How does an additive collection
\[
S\subseteq M
\]
produce an affine toric variety, a monomial parametrization, and a
system of binomial equations?
\end{guidingquestion}


\section{Affine semigroups, monomial maps, and equations}
\label{sec:affine-semigroups-monomial-maps}
\index{affine semigroup}
\index{semigroup algebra}
\index{monomial map}
\index{toric ideal}
\index{binomial equation}

Section~\ref{sec:affine-plane-to-affine-toric-varieties} suggested
that an affine toric variety should be determined by the characters
that extend regularly from its dense torus. Because products of
extending characters still extend, their lattice points form an
additive collection in the character lattice.

We now make that idea precise.

\subsection{Affine semigroups}

Throughout this chapter,
\[
\NN=\{0,1,2,\ldots\}.
\]

\begin{definition}
Let \(M\) be a lattice. An \emph{affine semigroup} is a finitely
generated additive subsemigroup
\[
S\subseteq M
\]
containing \(0\).

Equivalently, there are lattice points
\[
m_1,\ldots,m_r\in M
\]
such that
\[
S
=
\NN m_1+\cdots+\NN m_r.
\]
\end{definition}

We often write
\[
S=\NN\mathcal A
\qquad\text{for}\qquad
\mathcal A=\{m_1,\ldots,m_r\}.
\]

The subgroup of \(M\) generated by \(S\) will be denoted
\[
\ZZ S
=
\left\{
a_1m_1+\cdots+a_rm_r:
a_i\in\ZZ
\right\}.
\]
Although \(S\) remembers which nonnegative combinations are allowed,
\(\ZZ S\) remembers the lattice of all differences between elements
of \(S\).

\begin{example}
Inside \(M=\ZZ\), consider
\[
S=\NN\{2,3\}.
\]
Then
\[
S=\{0,2,3,4,5,\ldots\}.
\]
The integer \(1\) does not lie in \(S\), even though
\[
1=3-2\in\ZZ S.
\]
This missing lattice point will become important when we discuss
normality.
\end{example}

\begin{pauseandthink}
Find finite generating sets for the following affine semigroups:
\[
\ZZ_{\geq0}^2,
\qquad
2\ZZ_{\geq0},
\qquad
\{(a,b)\in\ZZ^2:a\geq0,\ a+b\geq0\}.
\]
A generating set need not be unique. Can one of your generating sets
be made smaller?
\end{pauseandthink}

\subsection{The semigroup algebra}

An affine semigroup can be turned directly into an algebra.

\begin{definition}
Let \(S\subseteq M\) be an affine semigroup. Its \emph{semigroup
algebra} is
\[
\Bbbk[S]
=
\bigoplus_{m\in S}\Bbbk\,\character{m},
\]
with multiplication determined by
\[
\character{m}\character{m'}
=
\character{m+m'}.
\]
\end{definition}

This notation treats the elements \(\character{m}\) as formal basis
vectors. Since \(S\subseteq M\), we may also regard them as actual
characters of the torus with character lattice \(M\). In that
realization,
\[
\Bbbk[S]\subseteq\Bbbk[M].
\]

If
\[
S=\NN m_1+\cdots+\NN m_r,
\]
then
\[
\Bbbk[S]
=
\Bbbk[
\character{m_1},\ldots,\character{m_r}
].
\]
Thus finite generation of the semigroup gives finite generation of
the algebra. The basic examples from the previous section become
\[
\begin{array}{c|c|c}
S&\Bbbk[S]&\Spec\Bbbk[S]\\
\hline
\NN^n&\Bbbk[x_1,\ldots,x_n]&\AAA^n\\
\ZZ_{\geq0}\times\ZZ&\Bbbk[x,y^{\pm1}]&\AAA^1\times\Gm\\
\ZZ^n&\Bbbk[t_1^{\pm1},\ldots,t_n^{\pm1}]&\Gm^n.
\end{array}
\]

The notation \(\Bbbk[S]\) encodes both the basis and the
multiplication. Adding lattice points corresponds to multiplying
monomials:
\[
\begin{array}{c@{\qquad\longleftrightarrow\qquad}c}
m+m'
&
\character{m}\character{m'}.
\end{array}
\]

\subsection{From generators to a monomial map}

Let
\[
S=\NN\mathcal A,
\qquad
\mathcal A=\{m_1,\ldots,m_r\}\subseteq M.
\]
The generators determine a monomial map
\[
\Phi_{\mathcal A}:T\longrightarrow\AAA^r
\]
defined by
\[
\Phi_{\mathcal A}(t)
=
\bigl(
\character{m_1}(t),\ldots,\character{m_r}(t)
\bigr).
\]

Every coordinate of the map is a character. We therefore obtain a
parameterized subset of affine space using monomials alone.

\begin{definition}
The \emph{affine semigroup variety} associated to
\(\mathcal A\) is the Zariski closure
\[
Y_{\mathcal A}
=
\overline{\Phi_{\mathcal A}(T)}
\subseteq\AAA^r.
\]
\end{definition}

The closure is necessary: the image of the torus has all coordinates
nonzero, while the closure adds the boundary points obtained when
some monomials approach zero.

\begin{example}
Take
\[
\mathcal A=\{(1,0),(1,1),(0,1)\}\subseteq\ZZ^2.
\]
The monomial map is
\[
\Phi_{\mathcal A}(s,t)
=
(s,st,t).
\]
Writing the coordinates of \(\AAA^3\) as \(x,z,y\), the image satisfies
\[
z=xy.
\]
Its closure is the surface
\[
Y_{\mathcal A}
=
V(z-xy)\subseteq\AAA^3,
\]
which is isomorphic to \(\AAA^2\). The extra generator \((1,1)\) did
not change the abstract variety; it produced a redundant affine
embedding.
\end{example}

\begin{proposition}
\label{prop:semigroup-algebra-monomial-closure}
Let
\[
S=\NN\mathcal A\subseteq M
\]
be an affine semigroup. Then:
\begin{enumerate}[label=(\alph*),wide=0pt]
\item \(\Bbbk[S]\) is a finitely generated integral domain;
\item \(\Spec\Bbbk[S]\) is an affine toric variety whose dense torus
has character lattice \(\ZZ S\);
\item there is a natural isomorphism
\[
\Spec\Bbbk[S]\simeq Y_{\mathcal A}.
\]
In particular,
\[
\dim\Spec\Bbbk[S]
=
\rank(\ZZ S).
\]
\end{enumerate}
\end{proposition}

\noindent
\textit{Reference.}
Over \(\CC\), this combines Cox--Little--Schenck,
Propositions~1.1.8 and~1.1.14. The same argument works over our
algebraically closed base field \(\Bbbk\).

\begin{proof}[Conceptual proof sketch]
The generators \(m_1,\ldots,m_r\) give a homomorphism
\[
\pi:
\Bbbk[x_1,\ldots,x_r]
\longrightarrow
\Bbbk[M],
\qquad
x_i\longmapsto\character{m_i}.
\]
Its image is exactly
\[
\Bbbk[
\character{m_1},\ldots,\character{m_r}
]
=
\Bbbk[S].
\]
On the geometric side, \(\pi\) is the pullback of the monomial map
\(\Phi_{\mathcal A}\). Therefore the coordinate ring of the closure
of its image is the image of \(\pi\), namely \(\Bbbk[S]\).

The image of the torus is again a torus. Its character lattice is the
subgroup generated by the characters used in the map, which is
\(\ZZ S\). The closure is irreducible because it is the closure of an
irreducible torus, and the torus action preserves the closure.

To complete every detail, one checks that the ideal of the closure is
\(\ker\pi\), verifies that the image torus is open in the closure, and
identifies its character lattice through pullback.
\end{proof}

\begin{warningbox}
The dense torus of \(\Spec\Bbbk[S]\) has character lattice
\[
\ZZ S,
\]
which may be smaller than the ambient lattice \(M\) in which \(S\) was
first written.

Thus the intrinsic lattice of the variety is not automatically the
largest lattice visible in the notation. Working in the correct
lattice will be essential when we discuss saturation and normality.
\end{warningbox}

\subsection{Relations among generators become equations}

The monomial map gives a parametrization. To find equations for its
closure, we look for additive relations among the generators
\[
m_1,\ldots,m_r.
\]

Define the lattice homomorphism
\[
\beta:\ZZ^r\longrightarrow M,
\qquad
e_i\longmapsto m_i,
\]
and let
\[
L=\ker\beta.
\]
An element
\[
\ell=(\ell_1,\ldots,\ell_r)\in L
\]
records an integral relation
\[
\ell_1m_1+\cdots+\ell_rm_r=0.
\]

Write
\[
\ell=\ell^+-\ell^-,
\]
where \(\ell^+,\ell^-\in\NN^r\) have disjoint supports. The relation
\(\beta(\ell)=0\) becomes
\[
\sum_i\ell_i^+m_i
=
\sum_i\ell_i^-m_i.
\]
Consequently,
\[
x^{\ell^+}-x^{\ell^-}
\]
vanishes after the substitution
\[
x_i=\character{m_i}.
\]

\begin{example}
For
\[
\mathcal A=\{(1,0),(1,1),(0,1)\},
\]
the relation
\[
(1,0)+(0,1)=(1,1)
\]
gives
\[
x_1x_3-x_2.
\]
This is the equation \(xy-z=0\) found above.
\end{example}

\begin{definition}
The kernel
\[
I_{\mathcal A}
=
\ker\left(
\Bbbk[x_1,\ldots,x_r]
\longrightarrow
\Bbbk[S]
\right),
\qquad
x_i\longmapsto\character{m_i},
\]
is the \emph{toric ideal} associated to \(\mathcal A\).
\end{definition}

\begin{proposition}[Equations of a monomial image]
\label{prop:toric-ideal-binomial-relations}
With notation as above,
\[
I_{\mathcal A}
=
\left\langle
x^{\ell^+}-x^{\ell^-}:
\ell\in L
\right\rangle.
\]
Equivalently, \(I_{\mathcal A}\) is generated by all binomials
\[
x^\alpha-x^\beta
\]
for which
\[
\sum_i\alpha_i m_i
=
\sum_i\beta_i m_i.
\]
\end{proposition}

\noindent
\textit{Reference.}
This is Cox--Little--Schenck, Proposition~1.1.9.

\begin{proof}[Conceptual proof sketch]
Every displayed binomial lies in the kernel because its two monomials
are sent to the same character.

For the reverse inclusion, take a polynomial relation among the
monomials \(\character{m_i}\). After substitution, different terms can
cancel only when they become the same character of the torus.
Group together all terms with the same resulting exponent in \(M\).
Within each group, subtract one chosen monomial from the others.
Every such difference is one of the binomials above.

A complete proof organizes this reduction using a monomial order, or
equivalently decomposes the polynomial according to its \(M\)-degree.
The driving idea is simply that distinct characters are linearly
independent in \(\Bbbk[M]\).
\end{proof}

Thus additive relations in the lattice become binomial equations in
affine space:
\[
\boxed{
\text{same sum of exponent vectors}
\quad\Longleftrightarrow\quad
\text{same monomial after parametrization}.
}
\]

\subsection{The cusp: the simplest missing exponent}

Let
\[
S=\NN\{2,3\}\subseteq\ZZ.
\]
Then
\[
\Bbbk[S]
=
\Bbbk[t^2,t^3].
\]
The monomial map is
\[
\Phi_{\mathcal A}:\Gm\longrightarrow\AAA^2,
\qquad
t\longmapsto(t^2,t^3).
\]

Let the coordinates of \(\AAA^2\) be \(x,y\), corresponding to the
generators \(2,3\). The relation
\[
3\cdot2=2\cdot3
\]
gives
\[
x^3-y^2.
\]
Hence
\[
\Bbbk[S]
\simeq
\Bbbk[x,y]/(x^3-y^2)
\]
and
\[
\Spec\Bbbk[S]
\simeq
V(x^3-y^2)\subseteq\AAA^2.
\]

This is the cuspidal cubic. It is an affine toric variety, even though
it is not normal. The character \(t=\character{1}\) is regular on the
normalization \(\AAA^1\), but it does not belong to
\(\Bbbk[t^2,t^3]\).

\begin{agcheatsheet}
For an irreducible affine variety \(X=\Spec A\), being
\emph{normal} means that \(A\) is integrally closed in its fraction
field. Concretely, if \(f\in\Frac(A)\) satisfies a monic
equation
\[
f^r+a_{r-1}f^{r-1}+\cdots+a_0=0,
\qquad
a_i\in A,
\]
then normality requires \(f\in A\).

Normal varieties may still be singular, but they are regular in
codimension one; consequently, their singularities can occur only in
codimension at least two. Every irreducible variety \(X\) has a
canonical normalization
\[
\nu:X^\nu\longrightarrow X,
\]
a finite birational morphism from a normal variety, unique up to
unique isomorphism.

For further reading, see Hartshorne, \emph{Algebraic Geometry},
Exercise~I.3.17; the Stacks Project, Tags~031O and~0BXQ; and
Cox--Little--Schenck, \S1.3.
\end{agcheatsheet}

\begin{pauseandthink}
Why does
\[
t\notin\Bbbk[t^2,t^3]
\]
even though
\[
t=\frac{t^3}{t^2}
\]
belongs to the fraction field? Why does
\[
t^2\in\Bbbk[t^2,t^3]
\]
suggest that the missing exponent \(1\) should be relevant to
normality?
\end{pauseandthink}

\subsection{A quadratic cone and the correct lattice}

Consider
\[
\mathcal A
=
\{(2,0),(1,1),(0,2)\}
\subseteq\ZZ^2
\]
and
\[
S=\NN\mathcal A.
\]
The monomial map
\[
\Gm^2\longrightarrow\AAA^3
\]
is
\[
(s,t)\longmapsto(s^2,st,t^2).
\]
Writing the coordinates of \(\AAA^3\) as \(a,b,c\), the relation
\[
(2,0)+(0,2)=2(1,1)
\]
gives
\[
ac-b^2=0.
\]
Therefore
\[
\Spec\Bbbk[S]
\simeq
V(ac-b^2)\subseteq\AAA^3.
\]

At first glance, \(S\) appears to omit many lattice points of the
first quadrant, such as \((1,0)\). But the character lattice of the
dense torus is not all of \(\ZZ^2\). It is
\[
\ZZ S
=
\{(a,b)\in\ZZ^2:a+b\text{ is even}\}.
\]
Inside this intrinsic lattice, the apparent gaps look very different.

This example anticipates an important rule:

\begin{center}
\fbox{\parbox{0.86\linewidth}{\centering
Questions about missing lattice points must be asked in the
character lattice of the dense torus.
}}
\end{center}

\subsection{Four equivalent descriptions}

We can now state the basic classification of arbitrary affine toric
varieties, before imposing normality.

\begin{theorem}[Equivalent constructions of affine toric varieties]
\label{thm:equivalent-affine-toric-constructions}
Let \(X\) be an affine variety. The following are equivalent:
\begin{enumerate}[label=(\alph*),wide=0pt]
\item \(X\) is an affine toric variety;
\item \(X\simeq Y_{\mathcal A}\) for a finite set
\(\mathcal A\) in a lattice;
\item \(X\) is isomorphic to an affine variety defined by a toric
ideal;
\item
\[
X\simeq\Spec\Bbbk[S]
\]
for an affine semigroup \(S\).
\end{enumerate}
\end{theorem}

\noindent
\textit{Reference.}
This is Cox--Little--Schenck, Theorem~1.1.17, adapted from \(\CC\) to
our algebraically closed field \(\Bbbk\).

\begin{proof}[Conceptual proof sketch]
Starting with a finite set of lattice points, the monomial map,
its binomial equations, and the semigroup algebra are three
presentations of the same object:
\[
Y_{\mathcal A}
\simeq
\Spec\Bbbk[\NN\mathcal A].
\]
The preceding propositions establish these implications.

The only genuinely new direction starts with an intrinsic affine
toric variety \(X\). Since the dense torus \(T\) is open and dense,
restriction embeds
\[
\Bbbk[X]\hookrightarrow\Bbbk[T]=\Bbbk[M].
\]
The extended torus action makes \(\Bbbk[X]\) stable under \(T\).
The weight-space decomposition then says that this subalgebra is
spanned by the characters it contains:
\[
\Bbbk[X]
=
\bigoplus_{\character{m}\in\Bbbk[X]}
\Bbbk\,\character{m}.
\]
The set
\[
S_X
=
\{m\in M:\character{m}\in\Bbbk[X]\}
\]
is closed under addition because \(\Bbbk[X]\) is an algebra. Finally,
finite generation of the coordinate ring implies that \(S_X\) is
finitely generated. Hence
\[
\Bbbk[X]=\Bbbk[S_X].
\]

To complete the proof in detail, one verifies that a finite collection
of algebra generators can be replaced by the finitely many characters
appearing in their weight decompositions.
\end{proof}

The theorem says that we are free to move between four viewpoints:
\[
\begin{array}{c}
\text{an intrinsic torus action},
\\
\text{a monomial parametrization},
\\
\text{binomial equations},
\\
\text{a semigroup algebra}.
\end{array}
\]

\begin{guidingquestion}
An affine semigroup \(S\) generates a real cone
\[
\Cone{S}\subseteq M_{\RR}.
\]
When can \(S\) be recovered as all lattice points of that cone, and
what geometric property does this recovery encode?
\end{guidingquestion}


\section{Cones, saturation, and normality}
\label{sec:cones-saturation-normality}
\index{cone!rational polyhedral}
\index{Gordan's lemma}
\index{Hilbert basis}
\index{saturation}
\index{normality}
\index{normalization}

An affine semigroup
\[
S=\NN m_1+\cdots+\NN m_r\subseteq M
\]
contains discrete information: it records exactly which lattice points
are allowed. If we permit nonnegative real coefficients instead, we
obtain a continuous geometric object:
\[
\Cone{S}
=
\RR_{\geq0}m_1+\cdots+\RR_{\geq0}m_r
\subseteq M_{\RR}.
\]

Passing from \(S\) to \(\Cone{S}\) remembers the directions in which
the semigroup grows, but it may forget holes. The main question of
this section is whether the original semigroup can be recovered from
its cone:
\[
S\stackrel{?}{=}\Cone{S}\cap M.
\]

The answer is controlled simultaneously by a combinatorial condition
on \(S\) and an algebraic condition on \(\Bbbk[S]\):
\[
\boxed{
S\text{ has no hidden lattice-point holes}
\quad\Longleftrightarrow\quad
\Bbbk[S]\text{ is normal}.
}
\]

\subsection{From a semigroup to a cone}

\begin{definition}
Let \(V\) be a finite-dimensional real vector space. A
\emph{polyhedral cone} is a subset of the form
\[
C
=
\RR_{\geq0}v_1+\cdots+\RR_{\geq0}v_r
\subseteq V
\]
for finitely many vectors \(v_1,\ldots,v_r\in V\).

When \(V=M_{\RR}\) for a lattice \(M\), the cone is
\emph{rational} if the generators may be chosen in \(M\).
\end{definition}

Thus the cone of every affine semigroup is a rational polyhedral
cone.

\begin{example}
Let
\[
S
=
\NN(2,0)+\NN(0,1)+\NN(1,1)
\subseteq\ZZ^2.
\]
Its cone is again the first quadrant:
\[
\Cone{S}
=
\RR_{\geq0}^2.
\]
However,
\[
(1,0)\notin S,
\qquad
2(1,0)=(2,0)\in S.
\]
Hence
\[
S\subsetneq\Cone{S}\cap\ZZ^2.
\]
The cone remembers the first quadrant but does not remember the
missing point \((1,0)\).
\end{example}

\begin{figure}[ht]
\centering
\begin{tikzpicture}[scale=0.95]
  \fill[PastelYellow] (0,0) rectangle (4.2,3.7);
  \draw[->,black!55] (-0.45,0)--(5.4,0);
  \draw[->,black!55] (0,-0.45)--(0,4.05);

  \foreach \a in {0,...,4}{
    \foreach \b in {0,...,3}{
      \fill[black!20] (\a,\b) circle (1.6pt);
    }
  }

  % Points generated by the semigroup in the displayed window.
  \foreach \a/\b in {
    0/0,2/0,4/0,
    0/1,1/1,2/1,3/1,4/1,
    0/2,1/2,2/2,3/2,4/2,
    0/3,1/3,2/3,3/3,4/3
  }{
    \fill[DictionaryBlue] (\a,\b) circle (2.45pt);
  }

  \draw[WarningRed,very thick] (1,0) circle (5.2pt);
  \node[WarningRed,below] at (1,-0.14)
    {\((1,0)\) missing};

  \draw[PredictionTeal,very thick,-{Latex[length=2.4mm]}]
    (4.75,2.65)--(3.65,2.65);
  \node[
    fill=white,
    draw=PredictionTeal,
    rounded corners,
    align=left,
    text width=3.15cm,
    anchor=west,
    inner sep=5pt
  ] at (4.82,2.65)
    {\(\Cone{S}=\RR_{\geq0}^2\).\\
     The cone contains every lattice point in the shaded first quadrant,\\
     but \(S\) still misses \((1,0)\).};
\end{tikzpicture}
\caption{The semigroup generated by \((2,0),(0,1),(1,1)\) has the
whole first quadrant as its cone, but it misses \((1,0)\).}
\label{fig:semigroup-hole-first-quadrant}
\end{figure}

\begin{pauseandthink}
For
\[
S=\NN\{2,3\}\subseteq\ZZ,
\]
determine \(\Cone{S}\), \(\Cone{S}\cap\ZZ\), and the lattice points
missing from \(S\). Which of the missing points remain missing after
multiplication by every positive integer?
\end{pauseandthink}

\subsection{A small amount of cone geometry}

We record the pieces of cone language that will be used throughout
the course.

The \emph{span} of a cone \(C\) is the smallest linear subspace
containing it:
\[
\Span_{\RR}(C)=C-C.
\]
Its dimension is the dimension of this span.

A cone is \emph{strongly convex}, or \emph{pointed}, if it contains no
nonzero linear subspace:
\[
C\cap(-C)=\{0\}.
\]
A cone that contains a line has directions in which one may move both
forward and backward while remaining inside the cone.

\begin{definition}
Let \(C\subseteq M_{\RR}\) be a polyhedral cone. A \emph{face} of
\(C\) is a subset
\[
F=C\cap H
\]
where \(H\) is a supporting hyperplane: \(C\) lies in one of the two
closed half-spaces determined by \(H\).

A one-dimensional face of a strongly convex cone is called a
\emph{ray}. A codimension-one face is called a \emph{facet}.
\end{definition}

For a rational ray \(\rho\), the lattice points
\[
\rho\cap M
\]
are generated by a unique primitive lattice point. We call it the
\emph{primitive ray generator}.

\begin{observation}
Faces will acquire algebraic and geometric meanings later:

\[
\begin{array}{c@{\qquad\longleftrightarrow\qquad}c}
\text{faces of a cone}
&
\text{localizations and torus orbits}.
\end{array}
\]

For the moment, faces simply describe the boundary structure of a
cone.
\end{observation}

\noindent
\textit{Reference.}
The basic cone terminology appears in Cox--Little--Schenck,
Definitions~1.2.2, 1.2.5, and~1.2.14. Primitive ray generators are
discussed immediately before Lemma~1.2.15.

\subsection{Gordan's lemma}

Starting with a finite set of lattice points always produces a
finitely generated semigroup. The converse direction is less obvious:
if we start with a rational cone and take \emph{all} of its lattice
points, do we still obtain a finitely generated semigroup?

\begin{theorem}[Gordan's lemma]
\label{thm:gordans-lemma-character-cone}
Let \(C\subseteq M_{\RR}\) be a rational polyhedral cone. Then
\[
C\cap M
\]
is a finitely generated semigroup.
\end{theorem}

\noindent
\textit{Reference.}
This is Cox--Little--Schenck, Proposition~1.2.17, applied to a
rational polyhedral cone in the character space \(M_{\RR}\).

\begin{proof}[Conceptual proof sketch]
Choose lattice generators
\[
v_1,\ldots,v_r\in M
\]
for \(C\). Every point of \(C\) is a nonnegative real combination
\[
a_1v_1+\cdots+a_rv_r.
\]
Split each coefficient into its integer and fractional parts:
\[
a_i=\lfloor a_i\rfloor+\delta_i,
\qquad
0\leq\delta_i<1.
\]
The integer parts produce a nonnegative integral combination of the
\(v_i\). The fractional remainder lies in the bounded region
\[
K=
\left\{
\delta_1v_1+\cdots+\delta_rv_r:
0\leq\delta_i<1
\right\}.
\]
Since a bounded region contains only finitely many lattice points,
\(K\cap M\) is finite. Every lattice point of \(C\) is therefore the
sum of one of these finitely many remainder points and a nonnegative
integral combination of the \(v_i\).

To complete the proof, retain from \(K\cap M\) only the points that
actually lie in \(C\), and add them to the original generators.
\end{proof}

The proof explains the finite generation geometrically:

\begin{center}
\fbox{\parbox{0.86\linewidth}{\centering
Integer parts move along the cone generators, while only finitely
many fractional remainders must be remembered.
}}
\end{center}

\subsection{The Hilbert basis}

When \(C\) is strongly convex, the semigroup \(C\cap M\) has a
canonical minimal generating set.

\begin{definition}
Let \(S=C\cap M\), where \(C\) is strongly convex. A nonzero element
\(m\in S\) is \emph{irreducible} if
\[
m=m'+m'',
\qquad
m',m''\in S,
\]
implies \(m'=0\) or \(m''=0\).

The set of irreducible elements of \(S\) is called the
\emph{Hilbert basis} of \(C\cap M\).
\end{definition}

\begin{proposition}
\label{prop:hilbert-basis-cone}
The Hilbert basis of \(C\cap M\) is finite, generates \(C\cap M\),
and is its unique minimal generating set.
\end{proposition}

\noindent
\textit{Reference.}
This is the character-space version of Cox--Little--Schenck,
Proposition~1.2.23.

\begin{proof}[Conceptual proof sketch]
Gordan's lemma supplies some finite generating set. Any irreducible
element must occur in every generating set, because it cannot be
written as a sum of smaller nonzero semigroup elements. Conversely,
repeatedly decompose any reducible element. Strong convexity provides
a linear function that is strictly positive away from the origin, so
this process cannot continue indefinitely. The resulting pieces are
irreducible.

Thus the irreducible elements generate and are forced to occur in
every generating set.
\end{proof}

\begin{warningbox}
The primitive ray generators of \(C\) need not generate the whole
semigroup \(C\cap M\). A cone may have very few rays but require
additional lattice points between them in its Hilbert basis.

Cone generators and semigroup generators answer different questions.
\end{warningbox}

\subsection{Saturation fills the holes}

We now give a name to the condition suggested by the cusp and by
\cref{fig:semigroup-hole-first-quadrant}.

\begin{definition}
Let \(S\) be an affine semigroup, and let
\[
M_S=\ZZ S
\]
be its group of differences. The semigroup \(S\) is
\emph{saturated} if
\[
m\in M_S,\quad k>0,\quad km\in S
\qquad\Longrightarrow\qquad
m\in S.
\]
\end{definition}

The use of the intrinsic lattice \(M_S=\ZZ S\) is essential. A point
outside \(M_S\) is not a character of the dense torus of
\(\Spec\Bbbk[S]\), so it should not count as a missing character.

\begin{example}
The semigroup
\[
S=\NN\{2,3\}\subseteq\ZZ
\]
is not saturated because
\[
1\notin S,
\qquad
2\cdot1=2\in S.
\]
\end{example}

\begin{example}
For
\[
S=\NN(2,0)+\NN(0,1)+\NN(1,1)\subseteq\ZZ^2,
\]
we have
\[
\ZZ S=\ZZ^2.
\]
The point \((1,0)\) witnesses the failure of saturation:
\[
(1,0)\notin S,
\qquad
2(1,0)\in S.
\]
\end{example}

\begin{proposition}[Saturation and lattice points]
\label{prop:saturation-cone-lattice-points}
Let \(S\) be an affine semigroup and set
\[
M_S=\ZZ S,
\qquad
C=\Cone{S}\subseteq(M_S)_{\RR}.
\]
Then
\[
S\text{ is saturated}
\quad\Longleftrightarrow\quad
S=C\cap M_S.
\]
\end{proposition}

\noindent
\textit{Reference.}
This is Cox--Little--Schenck, Exercise~1.3.4(a).

\begin{proof}
The inclusion
\[
S\subseteq C\cap M_S
\]
always holds.

Suppose first that \(S\) is saturated, and let \(m\in C\cap M_S\).
Because \(C\) is a rational cone generated by elements of \(S\), the
point \(m\) can be expressed using nonnegative rational coefficients:
\[
m=q_1m_1+\cdots+q_rm_r,
\qquad
q_i\in\QQ_{\geq0}.
\]
Choose a positive integer \(k\) clearing all denominators. Then
\[
km\in S.
\]
Saturation gives \(m\in S\).

Conversely, suppose
\[
S=C\cap M_S.
\]
If \(km\in S\) for \(m\in M_S\) and \(k>0\), then
\[
m=\frac1k(km)\in C.
\]
Hence
\[
m\in C\cap M_S=S.
\]
\end{proof}

This proposition is the precise version of the slogan:

\[
\boxed{
\text{saturation fills every lattice-point hole inside the cone.}
}
\]

\subsection{Normality is saturation in algebraic form}

Recall from the algebraic geometry cheatsheet in
Section~\ref{sec:affine-semigroups-monomial-maps} that an affine
domain \(A\) is normal when it is integrally closed in its fraction
field.

For semigroup algebras, integral closure has a direct combinatorial
meaning.

\begin{theorem}[Normality criterion for affine toric varieties]
\label{thm:normality-saturation-cone}
Let \(S\) be an affine semigroup, and let
\[
M_S=\ZZ S
\]
be its intrinsic lattice. Set
\[
C=\Cone{S}\subseteq (M_S)_{\RR}.
\]
The following are equivalent:
\begin{enumerate}[label=(\alph*),wide=0pt]
\item \(\Spec\Bbbk[S]\) is normal;
\item \(S\) is saturated in \(M_S\);
\item
\[
S=C\cap M_S.
\]
\end{enumerate}

Consequently, every normal affine toric variety arising from \(S\)
has coordinate ring
\[
\Bbbk[C\cap M_S],
\]
where \(C=\Cone{S}\subseteq(M_S)_{\RR}\). More intrinsically, if an
affine toric variety has dense-torus character lattice \(M_T\), then it
has the form
\[
X_C
=
\Spec\Bbbk[C\cap M_T]
\]
for a full-dimensional rational polyhedral cone
\(C\subseteq(M_T)_{\RR}\).
\end{theorem}

\noindent
\textit{Reference.}
This is the character-space formulation of
Cox--Little--Schenck, Theorem~1.3.5.

\begin{proof}[Conceptual proof sketch]
The implication
\[
\text{normal}\Longrightarrow\text{saturated}
\]
contains the key algebraic idea. Suppose
\[
km\in S
\]
for some \(m\in M_S\) and \(k>0\). Since \(m\in\ZZ S\), the character
\(\character{m}\) belongs to the fraction field of \(\Bbbk[S]\).
Moreover, it satisfies the monic equation
\[
Z^k-\character{km}=0
\]
with coefficient \(\character{km}\in\Bbbk[S]\). Thus
\(\character{m}\) is integral over \(\Bbbk[S]\). Normality forces
\(\character{m}\in\Bbbk[S]\), so \(m\in S\).

The equivalence between saturation and
\[
S=C\cap M_S
\]
is \cref{prop:saturation-cone-lattice-points}.

For the remaining direction, let
\[
N_S=\Hom_{\ZZ}(M_S,\ZZ)
\]
be the lattice dual to \(M_S\), and write
\[
\sigma=C^\vee\subseteq (N_S)_{\RR}.
\]
Since \(C\) spans \((M_S)_{\RR}\), the cone \(\sigma\) is strongly
convex.
The algebra \(\Bbbk[C\cap M_S]\) may be recovered as an intersection of
the algebras associated to the rays of \(\sigma\). After choosing a
basis adapted to a primitive ray generator, each ray algebra becomes
a polynomial ring in one variable with Laurent variables in the other
directions:
\[
\Bbbk[x_1,x_2^{\pm1},\ldots,x_n^{\pm1}].
\]
These rings are normal, and an intersection of integrally closed
domains inside a common fraction field is integrally closed.

The technical part of the complete proof verifies the intersection
description and the reduction to primitive ray generators.
\end{proof}

\begin{pauseandthink}
Return to
\[
S=\NN(2,0)+\NN(0,1)+\NN(1,1).
\]
Determine its saturation. What is the normalization of
\[
\Spec\Bbbk[S]?
\]
Using
\[
a=x^2,\qquad b=y,\qquad c=xy,
\]
find an equation for the original affine semigroup variety and write
the normalization map explicitly.
\end{pauseandthink}

Here
\[
\overline S=\ZZ_{\geq0}^2.
\]
The original algebra is
\[
\Bbbk[S]
=
\Bbbk[x^2,y,xy]
\simeq
\Bbbk[a,b,c]/(c^2-ab^2),
\]
while the saturated algebra is
\[
\Bbbk[\overline S]
=
\Bbbk[x,y].
\]
The normalization map is
\[
\AAA^2\longrightarrow V(c^2-ab^2)\subseteq\AAA^3,
\qquad
(x,y)\longmapsto(x^2,y,xy).
\]

\subsection{Normalization of an affine semigroup variety}

The previous theorem identifies the ring that should normalize
\(\Bbbk[S]\).

\begin{proposition}[Normalization by filling the cone]
\label{prop:normalization-semigroup-variety}
Let \(S\) be an affine semigroup, set
\[
M_S=\ZZ S,
\qquad
C=\Cone{S}\subseteq(M_S)_{\RR},
\]
and write
\[
\overline S=C\cap M_S.
\]
Then the inclusion
\[
\Bbbk[S]\subseteq\Bbbk[\overline S]
\]
induces the normalization morphism
\[
\nu:
\Spec\Bbbk[\overline S]
\longrightarrow
\Spec\Bbbk[S].
\]
\end{proposition}

\noindent
\textit{Reference.}
This is the character-space form of Cox--Little--Schenck,
Proposition~1.3.8.

\begin{proof}[Conceptual proof sketch]
By Gordan's lemma, \(\overline S=C\cap M_S\) is finitely generated.
It is saturated, so \(\Bbbk[\overline S]\) is normal.

For every \(m\in\overline S\), rationality of the cone implies that
some positive multiple \(km\) lies in \(S\). Hence
\[
\character{m}
\]
satisfies
\[
Z^k-\character{km}=0
\]
and is integral over \(\Bbbk[S]\). Thus
\(\Bbbk[\overline S]\) is integral, and therefore finite, over
\(\Bbbk[S]\).

Finally, \(S\) and \(\overline S\) generate the same group \(M_S\), so
their semigroup algebras have the same fraction field. The induced
morphism is therefore finite and birational, with normal source,
which characterizes the normalization.
\end{proof}

For the cusp this recovers the normalization already anticipated:
\[
\Bbbk[t^2,t^3]\subseteq\Bbbk[t],
\qquad
t\longmapsto(t^2,t^3).
\]

Different semigroups can generate the same cone, so cones do not
classify arbitrary affine toric varieties. If \(S_1\) and \(S_2\)
generate the same cone and the same intrinsic lattice, then they have
the same saturation and hence the same normalization.

All intersections with the cone must be taken in
\[
M_S=\ZZ S,
\]
the character lattice of the dense torus determined by \(S\). Using a
larger ambient lattice can create apparent holes that are not
intrinsic to the variety.

\begin{threedictionaries}
Normality is the same phenomenon in three languages.

\medskip

\begin{center}
\small
\begin{tabularx}{\linewidth}{@{}
  >{\raggedright\arraybackslash}X
  >{\raggedright\arraybackslash}X
  >{\raggedright\arraybackslash}X@{}}
\toprule
\textbf{Algebra}
&
\textbf{Geometry}
&
\textbf{Combinatorics}
\\
\midrule
\(\Bbbk[S]\)
&
\(\Spec\Bbbk[S]\)
&
an affine semigroup \(S\subseteq M\), with \(M_S=\ZZ S\)
\\[1mm]
integral closure
&
normalization
&
fill lattice-point holes
\\[1mm]
\(\Bbbk[S]\) integrally closed
&
the affine toric variety is normal
&
\(S\) is saturated
\\[1mm]
\(\Bbbk[C\cap M_S]\)
&
a normal affine toric variety
&
all lattice points of a rational cone \(C\)
\\
\bottomrule
\end{tabularx}
\end{center}
\end{threedictionaries}

\begin{guidingquestion}
If a cone
\[
C\subseteq M_{\RR}
\]
records the characters that extend regularly, what does the dual cone
in \(N_{\RR}\) record about one-parameter subgroups and their limits?
\end{guidingquestion}


\section{Dual cones and one-parameter subgroup limits}
\label{sec:dual-cones-oneps-limits}
\index{dual cone}
\index{one-parameter subgroup!limit criterion}
\index{strongly convex cone}

The previous section described a normal affine toric variety by a cone
\[
C\subseteq M_{\RR}
\]
whose lattice points index the characters that extend regularly. We
now turn this picture around.

Fix a torus \(T\), with character and one-parameter subgroup lattices
\[
M=X^*(T),
\qquad
N=X_*(T),
\]
and perfect pairing
\[
\pair{\phantom{m}}{\phantom{u}}:
M\times N\longrightarrow\ZZ.
\]

When a variety is first presented by an affine semigroup \(S\), the
lattices in this section should be understood intrinsically:
\[
M=M_S=\ZZ S,
\qquad
N=N_S=\Hom_{\ZZ}(M_S,\ZZ).
\]
An ambient lattice containing \(S\) plays no role in the geometry of
the dense torus unless it equals \(M_S\).

\subsection{Dual cones}

A vector \(u\in N_{\RR}\) defines a linear function on \(M_{\RR}\):
\[
m\longmapsto\pair{m}{u}.
\]
A cone in one of the two spaces therefore determines a cone in the
other by requiring all these values to be nonnegative.

\begin{definition}
Let
\[
\sigma\subseteq N_{\RR}
\]
be a polyhedral cone. Its \emph{dual cone} is
\[
\sigma^\vee
=
\left\{
m\in M_{\RR}:
\pair{m}{u}\geq0
\text{ for every }u\in\sigma
\right\}.
\]
Likewise, for a cone \(C\subseteq M_{\RR}\), its dual is
\[
C^\vee
=
\left\{
u\in N_{\RR}:
\pair{m}{u}\geq0
\text{ for every }m\in C
\right\}.
\]
\end{definition}

Thus \(\sigma^\vee\) consists of the linear functions on
\(N_{\RR}\) that are nonnegative throughout \(\sigma\).

\noindent
\textit{Reference.}
This is Cox--Little--Schenck, Definition~1.2.3.

If
\[
\sigma=\Cone{u_1,\ldots,u_r},
\]
then it is enough to test the generators:
\[
\sigma^\vee
=
\left\{
m\in M_{\RR}:
\pair{m}{u_i}\geq0
\text{ for }i=1,\ldots,r
\right\}.
\]
Indeed, nonnegativity on the \(u_i\) implies nonnegativity on every
nonnegative linear combination of them.

\begin{example}
Let
\[
\sigma
=
\Cone{(1,0),(1,2)}
\subseteq N_{\RR}.
\]
Then
\[
\sigma^\vee
=
\left\{
(a,b)\in M_{\RR}:
a\geq0,\ a+2b\geq0
\right\}.
\]
Its two rays are generated by
\[
(0,1)
\qquad\text{and}\qquad
(2,-1).
\]
\end{example}

\begin{figure}[ht]
\centering
\begin{tikzpicture}[scale=0.95]
  % N cone
  \begin{scope}[xshift=-4.6cm]
    \fill[SurfaceGreen!12]
      (0,0)--(3.1,0)--(1.55,3.1)--cycle;
    \draw[very thick,SurfaceGreen]
      (0,0)--(3.1,0);
    \draw[very thick,SurfaceGreen]
      (0,0)--(1.55,3.1);
    \draw[->,black!55] (-0.55,0)--(3.55,0);
    \draw[->,black!55] (0,-1.1)--(0,3.5);
    \draw[ray,->] (0,0)--(2.3,0);
    \draw[ray,->] (0,0)--(1.15,2.3);
    \node[below] at (2.35,0) {\((1,0)\)};
    \node[right] at (1.15,2.3) {\((1,2)\)};
    \node[font=\bfseries] at (1.1,3.75)
      {\(\sigma\subseteq N_{\RR}\)};
  \end{scope}

  % M dual cone
  \begin{scope}[xshift=4.3cm]
    \fill[DictionaryBlue!10]
      (0,0)--(0,3.1)--(3.2,-1.6)--cycle;
    \draw[very thick,DictionaryBlue]
      (0,0)--(0,3.1);
    \draw[very thick,DictionaryBlue]
      (0,0)--(3.2,-1.6);
    \draw[->,black!55] (-0.55,0)--(3.55,0);
    \draw[->,black!55] (0,-2.0)--(0,3.5);
    \draw[ray,->] (0,0)--(0,2.4);
    \draw[ray,->] (0,0)--(2.4,-1.2);
    \node[left] at (0,2.4) {\((0,1)\)};
    \node[above right] at (2.4,-1.2) {\((2,-1)\)};
    \node[font=\bfseries] at (1.1,3.75)
      {\(\sigma^\vee\subseteq M_{\RR}\)};
  \end{scope}

  \draw[very thick,<->,PredictionTeal]
    (-0.65,1.25)--(0.35,1.25);
  \node[PredictionTeal,above] at (-0.15,1.25)
    {dual};
\end{tikzpicture}
\caption{The cone
\(\sigma=\Cone{(1,0),(1,2)}\) and its dual
\(\sigma^\vee=\Cone{(0,1),(2,-1)}\).}
\label{fig:cone-and-dual-cone-example}
\end{figure}

\begin{proposition}[Biduality]
\label{prop:biduality-polyhedral-cones}
For every polyhedral cone,
\[
(\sigma^\vee)^\vee=\sigma.
\]
Moreover, the dual of a rational polyhedral cone is again rational
polyhedral.
\end{proposition}

\noindent
\textit{Reference.}
The biduality statement is Cox--Little--Schenck,
Proposition~1.2.4. Rationality of dual cones is recorded in the
discussion following Definition~1.2.14.

\begin{proof}[Conceptual proof sketch]
Every \(u\in\sigma\) pairs nonnegatively with every
\(m\in\sigma^\vee\), so
\[
\sigma\subseteq(\sigma^\vee)^\vee.
\]
For the reverse inclusion, a point outside a closed polyhedral cone
can be separated from it by a linear functional. That functional is
an element \(m\in\sigma^\vee\) that takes a negative value on the
outside point. Hence that point cannot lie in the bidual.

For a rational cone, the required supporting functionals may be
chosen rational and then cleared of denominators to become lattice
vectors.
\end{proof}

\begin{pauseandthink}
Compute the duals of the following cones in \(N_{\RR}\simeq\RR^2\):
\[
\{0\},
\qquad
\Cone{(1,0)},
\qquad
\Cone{(1,0),(0,1)},
\qquad
\Cone{(1,0),(-1,0)}.
\]
What happens to the dimension of the dual when the original cone
contains a line?
\end{pauseandthink}

\subsection{From a cone in \texorpdfstring{\(N_{\RR}\)}{N\_R} to an affine variety}

Let
\[
\sigma\subseteq N_{\RR}
\]
be a strongly convex rational polyhedral cone. Define
\[
S_\sigma
=
\sigma^\vee\cap M
\]
and
\[
U_\sigma
=
\Spec\Bbbk[S_\sigma]
=
\Spec\Bbbk[\sigma^\vee\cap M].
\]

Gordan's lemma makes \(S_\sigma\) finitely generated. Because it
contains every lattice point of the rational cone \(\sigma^\vee\), it
is saturated. Therefore \(U_\sigma\) is normal.

Strong convexity of \(\sigma\) ensures that
\[
\sigma^\vee
\]
spans \(M_{\RR}\). Consequently,
\[
\ZZ(\sigma^\vee\cap M)=M,
\]
so the dense torus of \(U_\sigma\) is the original torus \(T\).

\begin{theorem}
\label{thm:affine-toric-variety-from-N-cone}
Let \(\sigma\subseteq N_{\RR}\) be a strongly convex rational
polyhedral cone. Then
\[
U_\sigma
=
\Spec\Bbbk[\sigma^\vee\cap M]
\]
is a normal affine toric variety with dense torus \(T\).
\end{theorem}

\noindent
\textit{Reference.}
The construction and dense-torus statement are
Cox--Little--Schenck, Theorem~1.2.18 and Remark~1.2.19.
Normality follows from Theorem~1.3.5.

\begin{proof}[Conceptual proof sketch]
Gordan's lemma supplies finite generation, so \(U_\sigma\) is an
affine semigroup variety. The semigroup
\(\sigma^\vee\cap M\) is saturated, hence its algebra is normal.
Finally, strong convexity is dual to the condition that
\(\sigma^\vee\) be full-dimensional, which guarantees that its lattice
points generate the full character lattice \(M\).
\end{proof}

This construction translates the result of the previous section. If
\[
C\subseteq M_{\RR}
\]
is a full-dimensional rational polyhedral cone and
\[
X_C=\Spec\Bbbk[C\cap M],
\]
then
\[
\sigma=C^\vee\subseteq N_{\RR}
\]
is strongly convex and
\[
C=\sigma^\vee.
\]
Hence
\[
X_C=U_\sigma.
\]

The standard cones recover the basic affine toric varieties:
\[
\begin{array}{c|c|c}
\sigma&\sigma^\vee\cap M&U_\sigma\\
\hline
\{0\}&M&T\\
\Cone{(1,0)}\subseteq\RR^2
&\{(a,b)\in\ZZ^2:a\geq0\}
&\AAA^1\times\Gm\\
\Cone{(1,0),(0,1)}
&\ZZ_{\geq0}^2
&\AAA^2.
\end{array}
\]
The nonsmooth cone
\[
\Cone{(1,0),(1,2)}
\]
produces the normal quadratic cone
\[
V(AC-B^2)\subseteq\AAA^3,
\]
as computed in Section~\ref{sec:affine-semigroups-monomial-maps} and
used again in the smoothness discussion of
Section~\ref{sec:properties-affine-toric-varieties}.

\subsection{What does it mean for a limit to exist?}

For \(u\in N\), let
\[
\oneps{u}:\Gm\longrightarrow T
\]
be the corresponding one-parameter subgroup.

When we write
\[
\lim_{t\to0}\oneps{u}(t)
\quad\text{in }U_\sigma,
\]
we mean algebraically that the morphism \(\oneps{u}\) extends across
the missing parameter value:
\[
\begin{tikzcd}[column sep=large]
\Gm \arrow[r,"\oneps{u}"] \arrow[d,hook]
&
T \arrow[d,hook]
\\
\AAA^1 \arrow[r,dashed,"\overline{\oneps{u}}"']
&
U_\sigma.
\end{tikzcd}
\]
The limit is then the point
\[
\overline{\oneps{u}}(0).
\]

Because \(U_\sigma\) is affine, the extension can be tested on its
coordinate ring. For every
\[
m\in\sigma^\vee\cap M,
\]
we have
\[
(\oneps{u})^*(\character{m})
=
t^{\pair{m}{u}}.
\]
This Laurent monomial extends regularly to \(\AAA^1\) precisely when
\[
\pair{m}{u}\geq0.
\]

\subsection{The one-parameter subgroup limit criterion}

We have reached the geometric reason for placing the defining cone in
the one-parameter subgroup space.

\begin{theorem}[One-parameter subgroup limit criterion]
\label{thm:oneps-limit-criterion-affine}
Let
\[
\sigma\subseteq N_{\RR}
\]
be a strongly convex rational polyhedral cone, and let \(u\in N\).
Then
\[
u\in\sigma
\quad\Longleftrightarrow\quad
\oneps{u}:\Gm\longrightarrow T
\text{ extends to a morphism }
\AAA^1\longrightarrow U_\sigma.
\]
Equivalently,
\[
u\in\sigma
\quad\Longleftrightarrow\quad
\lim_{t\to0}\oneps{u}(t)
\text{ exists in }U_\sigma.
\]
\end{theorem}

\noindent
\textit{Reference.}
This is Cox--Little--Schenck, Proposition~3.2.2. Although it appears
there in the discussion of the orbit--cone correspondence, its proof
uses only the affine coordinate ring
\(\Bbbk[\sigma^\vee\cap M]\). The same algebraic proof works over our
base field \(\Bbbk\).

\begin{proof}
A morphism
\[
\AAA^1\longrightarrow U_\sigma
\]
is equivalent to a \(\Bbbk\)-algebra homomorphism
\[
\Bbbk[\sigma^\vee\cap M]\longrightarrow\Bbbk[t].
\]
The pullback along \(\oneps{u}\) is already defined with values in
\(\Bbbk[t,t^{-1}]\), and sends
\[
\character{m}\longmapsto t^{\pair{m}{u}}.
\]
It takes values in \(\Bbbk[t]\) exactly when
\[
\pair{m}{u}\geq0
\qquad
\text{for every }m\in\sigma^\vee\cap M.
\]
Since the lattice points generate the rational cone
\(\sigma^\vee\), this is equivalent to
\[
u\in(\sigma^\vee)^\vee.
\]
By biduality,
\[
(\sigma^\vee)^\vee=\sigma.
\]
\end{proof}

The theorem gives the promised second description of \(U_\sigma\):

\[
\boxed{
\sigma\cap N
=
\left\{
u\in N:
\oneps{u}\text{ has a limit in }U_\sigma
\text{ as }t\to0
\right\}.
}
\]

On the character side,
\[
\sigma^\vee\cap M
=
\left\{
m\in M:
\character{m}\text{ extends regularly to }U_\sigma
\right\}.
\]

Thus the two cones encode two complementary extension problems:
\[
\begin{array}{c|c}
\sigma^\vee\subseteq M_{\RR}
&
\sigma\subseteq N_{\RR}
\\
\hline
\text{which functions extend}
&
\text{which motions acquire limits}.
\end{array}
\]

\subsection{The limit point}

Suppose \(u\in\sigma\cap N\). Evaluating the extended map at \(t=0\)
gives the semigroup homomorphism
\[
\gamma_u:\sigma^\vee\cap M\longrightarrow\Bbbk
\]
defined by
\[
\gamma_u(m)
=
\begin{cases}
1,&\pair{m}{u}=0,\\
0,&\pair{m}{u}>0.
\end{cases}
\]
This is the point
\[
\lim_{t\to0}\oneps{u}(t)\in U_\sigma.
\]

If \(u\) lies in the relative interior of \(\sigma\), then
\[
\pair{m}{u}=0
\quad\Longleftrightarrow\quad
m\in\sigma^\perp\cap M.
\]
The resulting point is called the \emph{distinguished point}
\(\gamma_\sigma\) of \(U_\sigma\).

More generally, let \(\tau\preceq\sigma\) be the smallest face
containing \(u\). Then
\[
u\in\relint(\tau),
\]
and the limit depends only on \(\tau\):
\[
\gamma_u(m)
=
\begin{cases}
1,&m\in\tau^\perp\cap M,\\
0,&m\notin\tau^\perp.
\end{cases}
\]
We will later interpret this point as the distinguished point
associated to the face \(\tau\).

\begin{example}
Let
\[
U_\sigma=\AAA^2,
\qquad
\sigma=\RR_{\geq0}^2.
\]
For
\[
u=(a,b)\in\sigma\cap\ZZ^2,
\]
the one-parameter subgroup is
\[
\oneps{u}(t)=(t^a,t^b).
\]
Its limit is
\[
\begin{array}{c|c}
(a,b)
&
\displaystyle\lim_{t\to0}(t^a,t^b)
\\
\hline
a>0,\ b>0&(0,0),\\
a>0,\ b=0&(0,1),\\
a=0,\ b>0&(1,0),\\
a=b=0&(1,1).
\end{array}
\]
The four possibilities correspond to the relative interiors of the
four faces of the first quadrant: the two-dimensional cone, its two
rays, and the zero face.
\end{example}

\begin{surfacelaboratory}
For \(\AAA^2\), a one-parameter subgroup by itself reaches one
distinguished point on each orbit. Translating the motion by a torus
point,
\[
t\longmapsto p\,\oneps{u}(t),
\]
moves that distinguished limit through the corresponding orbit.

For example, if \(u=(1,0)\), then
\[
(a,b)\,(t,1)=(at,b)
\longrightarrow(0,b).
\]
As \(b\) varies through \(\Bbbk^\times\), these limits fill the
one-dimensional orbit
\[
V(x)\setminus\{(0,0)\}.
\]
\end{surfacelaboratory}

\begin{pauseandthink}
For
\[
\sigma=\Cone{(1,0),(1,2)}
\]
and \(u=(p,q)\in N=\ZZ^2\), use the inequalities defining \(\sigma\)
to determine exactly when
\[
\lim_{t\to0}\oneps{u}(t)
\]
exists in
\[
U_\sigma\simeq V(AC-B^2).
\]
Then compute the pullbacks of \(A,B,C\) and determine the limit point
when \(u\) lies:
\[
\text{(a) in }\relint(\sigma),\qquad
\text{(b) on each ray of }\sigma.
\]
\end{pauseandthink}

For a single affine variety, the cones
\[
\sigma\subseteq N_{\RR}
\qquad\text{and}\qquad
\sigma^\vee\subseteq M_{\RR}
\]
contain the same information. The advantage of \(\sigma\) is
geometric: its lattice points are one-parameter motions, and its faces
organize their limit points. This is why cones in \(N_{\RR}\) become
the natural objects for gluing affine toric varieties.

\begin{guidingquestion}
Which geometric properties of \(U_\sigma\)—its units, torus factors,
fixed points, smoothness, morphisms, and affine open subsets—can be
read directly from the cone \(\sigma\)?
\end{guidingquestion}


\section{Properties of affine toric varieties}
\label{sec:properties-affine-toric-varieties}
\index{semigroup homomorphism}
\index{torus factor}
\index{fixed point!of an affine toric variety}
\index{smoothness criterion!affine cone}
\index{toric morphism}
\index{face localization}

We now know how to construct a normal affine toric variety from a
strongly convex rational polyhedral cone:
\[
U_\sigma
=
\Spec\Bbbk[\sigma^\vee\cap M].
\]
The purpose of this section is to learn how much geometry can be read
directly from \(\sigma\).

Normality has already provided the first instance of this principle:
\[
U_\sigma\text{ is normal}
\quad\Longleftrightarrow\quad
\sigma^\vee\cap M\text{ contains every lattice point of }\sigma^\vee.
\]
We now add several more entries to the dictionary.

\subsection{Points as semigroup homomorphisms}

Let
\[
S\subseteq M
\]
be an affine semigroup. We regard \(\Bbbk\) as a multiplicative
semigroup, including the absorbing element \(0\).

\begin{proposition}[Points of an affine semigroup variety]
\label{prop:points-semigroup-homomorphisms}
For
\[
X_S=\Spec\Bbbk[S],
\]
there are natural correspondences among:
\begin{enumerate}[label=(\alph*),wide=0pt]
\item points \(p\in X_S\);
\item maximal ideals of \(\Bbbk[S]\);
\item semigroup homomorphisms
\[
\gamma:S\longrightarrow\Bbbk.
\]
\end{enumerate}
The homomorphism associated to \(p\) is
\[
\gamma_p(m)=\character{m}(p).
\]
\end{proposition}

\noindent
\textit{Reference.}
This is Cox--Little--Schenck, Proposition~1.3.1.

\begin{proof}
A point \(p\) gives the evaluation homomorphism
\[
\operatorname{ev}_p:\Bbbk[S]\longrightarrow\Bbbk.
\]
Restricting it to the monomial basis gives
\[
\gamma_p(m)=\operatorname{ev}_p(\character{m}).
\]
Since
\[
\character{m+m'}
=
\character{m}\character{m'},
\]
the map \(\gamma_p\) is multiplicative.

Conversely, a semigroup homomorphism \(\gamma:S\to\Bbbk\) extends
linearly to
\[
\Bbbk[S]\longrightarrow\Bbbk,
\qquad
\sum_m c_m\character{m}
\longmapsto
\sum_m c_m\gamma(m).
\]
Multiplicativity of \(\gamma\) makes this a \(\Bbbk\)-algebra
homomorphism, hence a point of \(\Spec\Bbbk[S]\).
\end{proof}

This description is especially effective because zeros are allowed.
A value
\[
\gamma(m)=0
\]
records that the character \(\character{m}\) vanishes at the point.

\begin{example}
For
\[
\AAA^2
=
\Spec\Bbbk[\NN^2],
\]
a semigroup homomorphism
\[
\gamma:\NN^2\longrightarrow\Bbbk
\]
is determined by
\[
a=\gamma(1,0),
\qquad
b=\gamma(0,1).
\]
Indeed,
\[
\gamma(i,j)=a^ib^j.
\]
Thus semigroup homomorphisms \(\NN^2\to\Bbbk\) are exactly the points
\[
(a,b)\in\AAA^2.
\]
\end{example}

The torus action also has an intrinsic semigroup description. If
\(t\in T\) and \(p\leftrightarrow\gamma_p\), then
\[
t\cdot p
\quad\longleftrightarrow\quad
\bigl(m\longmapsto\character{m}(t)\gamma_p(m)\bigr).
\]
The nonzero factors \(\character{m}(t)\) rescale the values of
\(\gamma_p\), but they do not change which values vanish.

\begin{pauseandthink}
A point \(p\in X_S\) belongs to the dense torus precisely when
\[
\gamma_p(m)\neq0
\qquad
\text{for every }m\in S.
\]
Explain why. Then identify the dense torus inside \(\AAA^2\) using
semigroup homomorphisms \(\NN^2\to\Bbbk\).
\end{pauseandthink}

\subsection{Invertible characters and torus factors}

A character \(\character{m}\) is invertible in
\[
\Bbbk[\sigma^\vee\cap M]
\]
exactly when both \(m\) and \(-m\) lie in \(\sigma^\vee\). This means
\[
\pair{m}{u}=0
\qquad
\text{for every }u\in\sigma.
\]

Define
\[
\sigma^\perp
=
\left\{
m\in M_{\RR}:
\pair{m}{u}=0
\text{ for every }u\in\sigma
\right\}.
\]

\begin{proposition}[Units]
\label{prop:units-affine-toric-coordinate-ring}
The units of the coordinate ring of \(U_\sigma\) are
\[
\Bbbk[U_\sigma]^\times
=
\left\{
c\,\character{m}:
c\in\Bbbk^\times,\ 
m\in\sigma^\perp\cap M
\right\}.
\]
\end{proposition}

\begin{proof}[Conceptual proof sketch]
Every displayed element is visibly invertible, with inverse
\[
c^{-1}\character{-m}.
\]

Conversely, suppose \(fg=1\). Pass from \(M\) to the quotient
\[
M/(\sigma^\perp\cap M).
\]
The image of \(\sigma^\vee\cap M\) in this quotient is pointed. Choose
a linear functional that is strictly positive on every nonzero
degree. Looking at the largest degree occurring in \(fg\) shows that
both \(f\) and \(g\) must have only one degree in the quotient.
Their remaining factors lie in the group algebra
\(\Bbbk[\sigma^\perp\cap M]\), whose units are scalar multiples of
characters.
\end{proof}

The orthogonal space \(\sigma^\perp\) therefore measures the
directions in which the affine toric variety still behaves like a
torus.

Let
\[
N_\sigma
=
N\cap\Span_{\RR}(\sigma).
\]
This is a saturated sublattice of \(N\) of rank
\[
r=\dim\sigma.
\]
Choose a complementary lattice \(N'\) such that
\[
N=N_\sigma\oplus N'.
\]
The choice is generally not canonical.

\begin{proposition}[Splitting off the torus factor]
\label{prop:affine-toric-torus-factor}
With notation as above,
\[
U_{\sigma,N}
\simeq
U_{\sigma,N_\sigma}\times T_{N'}.
\]
In particular,
\[
U_{\sigma,N}
\simeq
U_{\sigma,N_\sigma}\times\Gm^{\,\rank N-\dim\sigma}.
\]
\end{proposition}

\noindent
\textit{Reference.}
This is the decomposition in Cox--Little--Schenck,
Equation~(1.3.2), in the discussion following Theorem~1.3.12.

\begin{proof}[Conceptual proof sketch]
Dualizing the decomposition
\[
N=N_\sigma\oplus N'
\]
gives
\[
M=M_\sigma\oplus M'.
\]
Since \(\sigma\) lies entirely in \((N_\sigma)_{\RR}\), its dual
semigroup splits as
\[
\sigma^\vee\cap M
=
\bigl(\sigma^\vee\cap M_\sigma\bigr)\oplus M'.
\]
The corresponding semigroup algebra is therefore a tensor product:
\[
\Bbbk[\sigma^\vee\cap M]
\simeq
\Bbbk[\sigma^\vee\cap M_\sigma]\otimes_\Bbbk\Bbbk[M'].
\]
Taking spectra turns the tensor product into the displayed product of
varieties.
\end{proof}

\begin{example}
For the ray
\[
\sigma=\Cone{(1,0)}\subseteq\RR^2,
\]
we have
\[
U_\sigma\simeq\AAA^1\times\Gm.
\]
The character \(y\) lies in
\(\sigma^\perp\cap M=\ZZ(0,1)\), so \(y\) is invertible. The
\(\Gm\)-factor is visible both algebraically and geometrically.
\end{example}

\begin{threedictionaries}
The directions perpendicular to \(\sigma\) describe the same feature
in three languages.

\medskip

\begin{center}
\small
\begin{tabularx}{\linewidth}{@{}
  >{\raggedright\arraybackslash}X
  >{\raggedright\arraybackslash}X
  >{\raggedright\arraybackslash}X@{}}
\toprule
\textbf{Algebra}
&
\textbf{Geometry}
&
\textbf{Combinatorics}
\\
\midrule
invertible characters
\(\character{m}\)
&
torus directions that never meet a boundary
&
\(m\in\sigma^\perp\cap M\)
\\[1mm]
a Laurent polynomial factor
\(\Bbbk[M']\)
&
a factor \(T_{N'}\)
&
a complementary lattice to
\(N_\sigma\)
\\
\bottomrule
\end{tabularx}
\end{center}
\end{threedictionaries}

\subsection{Fixed points}

A torus-fixed point must remain unchanged under every rescaling by
the torus. In semigroup language this leaves only one possible
candidate.

\begin{definition}
An affine semigroup \(S\) is \emph{pointed} if
\[
S\cap(-S)=\{0\}.
\]
Equivalently, \(0\) is the only invertible element of \(S\).
\end{definition}

\begin{proposition}[Fixed points]
\label{prop:fixed-points-affine-toric}
Let
\[
X_S=\Spec\Bbbk[S].
\]
The torus action on \(X_S\) has a fixed point if and only if \(S\) is
pointed. When it exists, the fixed point is unique and is represented
by
\[
\gamma_0(m)
=
\begin{cases}
1,&m=0,\\
0,&m\neq0.
\end{cases}
\]
\end{proposition}

\noindent
\textit{Reference.}
This is Cox--Little--Schenck, Proposition~1.3.2.

\begin{proof}
If \(p\leftrightarrow\gamma\) is fixed, then
\[
\character{m}(t)\gamma(m)=\gamma(m)
\]
for every \(t\in T\). For \(m\neq0\), one can choose \(t\) with
\(\character{m}(t)\neq1\), so \(\gamma(m)=0\). Hence a fixed point, if
it exists, must be \(\gamma_0\).

The map \(\gamma_0\) is multiplicative exactly when two nonzero
elements of \(S\) cannot add to \(0\). This is equivalent to
\[
S\cap(-S)=\{0\}.
\]
\end{proof}

For the cone construction, pointedness has a simple geometric
translation.

\begin{corollary}
\label{cor:fixed-point-full-dimensional-cone}
Let
\[
\sigma\subseteq N_{\RR}
\]
be strongly convex. The torus action on \(U_\sigma\) has a fixed point
if and only if
\[
\dim\sigma=\dim N_{\RR}.
\]
When it exists, it is the distinguished point
\(\gamma_\sigma\) obtained as the limit of every one-parameter
subgroup in \(\relint(\sigma)\).
\end{corollary}

\noindent
\textit{Reference.}
The fixed-point criterion is Cox--Little--Schenck,
Corollary~1.3.3. The interpretation by one-parameter limits follows
from \cref{thm:oneps-limit-criterion-affine}.

\begin{proof}
The semigroup
\[
\sigma^\vee\cap M
\]
is pointed precisely when
\[
\sigma^\perp=\{0\}.
\]
This is equivalent to \(\sigma\) spanning \(N_{\RR}\), or
\[
\dim\sigma=\dim N_{\RR}.
\]
\end{proof}

\begin{example}
The affine plane corresponds to the full-dimensional first quadrant,
so it has one fixed point: the origin.

The variety
\[
\AAA^1\times\Gm
\]
corresponds to a ray in \(\RR^2\). It has no fixed point because the
\(\Gm\)-factor acts freely in its own direction.
\end{example}

\subsection{Smoothness}

The singularities of \(U_\sigma\) are controlled by how the primitive
ray generators of \(\sigma\) sit in the lattice.

\begin{definition}
A rational polyhedral cone
\[
\sigma\subseteq N_{\RR}
\]
is \emph{smooth} if its primitive ray generators form part of a basis
of \(N\).
\end{definition}

For a full-dimensional cone, smoothness means that its primitive ray
generators themselves form a basis of \(N\).

\begin{theorem}[Smoothness criterion]
\label{thm:smoothness-criterion-affine-toric}
Let
\[
\sigma\subseteq N_{\RR}
\]
be a strongly convex rational polyhedral cone. Then
\[
U_\sigma\text{ is smooth}
\quad\Longleftrightarrow\quad
\sigma\text{ is smooth}.
\]
Consequently, every smooth affine toric variety is isomorphic to
\[
\AAA^r\times\Gm^{n-r}
\]
for suitable \(r\leq n\).
\end{theorem}

\noindent
\textit{Reference.}
This is Cox--Little--Schenck, Theorem~1.3.12, together with the product
decomposition in Equation~(1.3.2).

\begin{proof}[Conceptual proof sketch]
If the primitive ray generators extend to a basis of \(N\), choose
that basis. In the dual basis,
\[
\sigma^\vee\cap M
\simeq
\NN^r\oplus\ZZ^{n-r}.
\]
Hence
\[
U_\sigma
\simeq
\AAA^r\times\Gm^{n-r},
\]
which is smooth.

For the converse, first split off the torus factor and reduce to a
full-dimensional cone. Its distinguished fixed point has maximal
ideal generated by the monomials corresponding to the Hilbert basis
of \(\sigma^\vee\cap M\). Modulo the square of that ideal, those
generators form a basis of the Zariski cotangent space. Smoothness
forces the tangent-space dimension to equal \(\dim U_\sigma=n\).

But every ray of \(\sigma^\vee\) contributes at least one Hilbert
basis element, and a full-dimensional cone has at least \(n\) rays.
Equality of dimensions forces exactly \(n\) rays and no additional
Hilbert basis elements. Those primitive ray generators must then form
a lattice basis. Dualizing shows that \(\sigma\) is smooth.

The complete proof carefully identifies the tangent space and handles
the reduction from a non-full-dimensional cone.
\end{proof}

\begin{example}
The first quadrant
\[
\Cone{(1,0),(0,1)}
\]
is smooth, and its affine toric variety is \(\AAA^2\).
\end{example}

\begin{example}
The cone
\[
\sigma=\Cone{(1,0),(1,2)}
\]
is not smooth because
\[
\det
\begin{pmatrix}
1&1\\
0&2
\end{pmatrix}
=2.
\]
Its primitive ray generators span a sublattice of index \(2\), rather
than a basis of \(\ZZ^2\). Correspondingly,
\[
U_\sigma\simeq V(AC-B^2)
\]
is singular at the origin.
\end{example}

\begin{surfacelaboratory}
For a two-dimensional full-dimensional cone
\[
\sigma=\Cone{u_1,u_2},
\]
smoothness is equivalent to
\[
|\det(u_1,u_2)|=1.
\]
Thus the determinant measures how far the two primitive ray
generators are from being a lattice basis.

The value \(1\) gives a smooth affine plane chart. A value larger than
\(1\) gives a toric surface singularity.
\end{surfacelaboratory}

\begin{pauseandthink}
Determine which of the following cones in \(\ZZ^2\) are smooth:
\[
\Cone{(1,0),(0,1)},\qquad
\Cone{(1,0),(1,3)},\qquad
\Cone{(1,1),(-1,0)}.
\]
For each smooth cone, write the associated affine toric variety as
\(\AAA^r\times\Gm^{2-r}\).
\end{pauseandthink}

\subsection{Toric morphisms}

The morphisms compatible with torus actions are also controlled by
lattice maps.

\begin{definition}
Let
\[
X_{S_1}=\Spec\Bbbk[S_1],
\qquad
X_{S_2}=\Spec\Bbbk[S_2].
\]
A morphism
\[
\phi:X_{S_1}\longrightarrow X_{S_2}
\]
is \emph{toric} if its pullback
\[
\phi^*:\Bbbk[S_2]\longrightarrow\Bbbk[S_1]
\]
is induced by a semigroup homomorphism
\[
\widehat\phi:S_2\longrightarrow S_1.
\]
\end{definition}

\noindent
\textit{Reference.}
This is Cox--Little--Schenck, Definition~1.3.13.

Such a morphism restricts to a group homomorphism between the dense
tori and is equivariant:
\[
\phi(t\cdot p)=\phi(t)\cdot\phi(p).
\]
Conversely, a morphism that maps the dense torus into the dense torus,
restricts there to a group homomorphism, and respects the torus
actions is toric.

\noindent
\textit{Reference.}
These statements are Cox--Little--Schenck,
Proposition~1.3.14.

For normal affine toric varieties, the criterion becomes particularly
simple.

\begin{proposition}[Cone criterion for toric morphisms]
\label{prop:cone-criterion-toric-morphism}
Let
\[
\sigma_i\subseteq(N_i)_{\RR},
\qquad
i=1,2,
\]
be strongly convex rational polyhedral cones, and let
\[
F:N_1\longrightarrow N_2
\]
be a lattice homomorphism. The induced torus homomorphism
\[
T_{N_1}\longrightarrow T_{N_2}
\]
extends to a toric morphism
\[
U_{\sigma_1}\longrightarrow U_{\sigma_2}
\]
if and only if
\[
F_{\RR}(\sigma_1)\subseteq\sigma_2.
\]
\end{proposition}

\noindent
\textit{Reference.}
This is Cox--Little--Schenck, Proposition~1.3.15.

\begin{proof}
Let
\[
F^*:M_2\longrightarrow M_1
\]
be the dual map. The torus homomorphism extends precisely when the
pullback of every regular character on \(U_{\sigma_2}\) remains
regular on \(U_{\sigma_1}\):
\[
F^*(\sigma_2^\vee\cap M_2)
\subseteq
\sigma_1^\vee\cap M_1.
\]

Suppose first that
\[
F_{\RR}(\sigma_1)\subseteq\sigma_2.
\]
For \(m\in\sigma_2^\vee\) and \(u\in\sigma_1\),
\[
\pair{F^*(m)}{u}
=
\pair{m}{F(u)}
\geq0.
\]
Thus
\[
F^*(m)\in\sigma_1^\vee.
\]

Conversely, suppose the pullback condition holds. By rationality, the
lattice points of \(\sigma_2^\vee\) generate the cone. Hence for every
\(u\in\sigma_1\) and every \(m\in\sigma_2^\vee\),
\[
\pair{m}{F(u)}
=
\pair{F^*(m)}{u}
\geq0.
\]
Biduality gives
\[
F(u)\in(\sigma_2^\vee)^\vee=\sigma_2.
\]
\end{proof}

\begin{example}
Let
\[
F:\ZZ^2\longrightarrow\ZZ,
\qquad
(a,b)\longmapsto a+b.
\]
The first quadrant maps into the ray
\[
\RR_{\geq0}\subseteq\RR.
\]
Therefore \(F\) induces the toric morphism
\[
\AAA^2\longrightarrow\AAA^1,
\qquad
(x,y)\longmapsto xy.
\]
Indeed, the target character \(z\) pulls back to \(xy\).
\end{example}

\begin{prediction}
Suppose two cones share a common face. The identity lattice map sends
that face into each cone. What morphisms should this produce between
the corresponding affine toric varieties, and what kind of subsets
should their images be?
\end{prediction}

\subsection{Faces give affine open subsets}

Let
\[
\tau\preceq\sigma
\]
be a face. There is a lattice point
\[
m\in\sigma^\vee\cap M
\]
such that
\[
\tau
=
\sigma\cap m^\perp
=
\{u\in\sigma:\pair{m}{u}=0\}.
\]
The character \(\character{m}\) is nonnegative along all directions
of \(\sigma\), and it vanishes in order precisely zero along the
directions of \(\tau\).

\begin{proposition}[Face localization]
\label{prop:face-localization-affine-toric}
With notation as above,
\[
\Bbbk[\tau^\vee\cap M]
=
\Bbbk[\sigma^\vee\cap M]_{\character{m}}.
\]
Consequently,
\[
U_\tau
\simeq
(U_\sigma)_{\character{m}}
=
D(\character{m})
\subseteq U_\sigma.
\]
Thus every face determines a torus-invariant affine open subset.
\end{proposition}

\noindent
\textit{Reference.}
This is Cox--Little--Schenck, Proposition~1.3.16 and
Equation~(1.3.4).

\begin{proof}[Conceptual proof sketch]
Localizing at \(\character{m}\) allows arbitrary negative powers of
that character. On exponent vectors, this replaces the semigroup
\[
\sigma^\vee\cap M
\]
by
\[
(\sigma^\vee\cap M)+\ZZ(-m).
\]
The defining property
\[
\tau=\sigma\cap m^\perp
\]
implies that this enlarged semigroup is exactly
\[
\tau^\vee\cap M.
\]
Taking spectra identifies the localization with the principal open
subset where \(\character{m}\neq0\).

The only substantive verification is the equality of the two
semigroups; geometrically, subtracting multiples of \(m\) removes the
inequality transverse to the face while retaining all inequalities
along \(\tau\).
\end{proof}

\begin{example}
Let
\[
\sigma=\Cone{e_1,e_2}\subseteq\RR^2,
\qquad
U_\sigma=\AAA^2.
\]
For the face
\[
\tau=\Cone{e_1},
\]
take
\[
m=e_2^\vee.
\]
Then
\[
\character{m}=y,
\]
and
\[
U_\tau
=
D(y)
=
\Spec\Bbbk[x,y^{\pm1}]
\simeq
\AAA^1\times\Gm.
\]
\end{example}

\begin{example}
For the zero face
\[
\tau=\{0\},
\]
we may take
\[
m=e_1^\vee+e_2^\vee.
\]
Then
\[
D(\character{m})
=
D(xy)
=
\Gm^2,
\]
the dense torus of \(\AAA^2\).
\end{example}

\begin{figure}[ht]
\centering
\begin{tikzpicture}[scale=1.05]
  % cone side
  \begin{scope}[xshift=-4.3cm]
    \fill[SurfaceGreen!10] (0,0)--(3.0,0)--(0,3.0)--cycle;
    \draw[very thick,SurfaceGreen] (0,0)--(3.0,0);
    \draw[very thick,SurfaceGreen] (0,0)--(0,3.0);
    \draw[very thick,WarningRed] (0,0)--(3.0,0);
    \node[font=\bfseries] at (1.25,3.5)
      {\(\tau\preceq\sigma\)};
    \node[WarningRed,below] at (1.8,0)
      {\(\tau=\Cone{e_1}\)};
    \node at (1.2,1.25) {\(\sigma\)};
  \end{scope}

  \draw[very thick,-{Latex[length=2.5mm]},PredictionTeal]
    (-0.6,1.35)--(0.55,1.35);
  \node[PredictionTeal,above] at (0,1.35)
    {localize};

  % variety side
  \begin{scope}[xshift=4.2cm]
    \draw[->,black!55] (-2.7,0)--(2.85,0);
    \draw[->,black!55] (0,-2.4)--(0,2.55);
    \draw[very thick,DictionaryBlue] (-2.5,0)--(2.5,0);
    \draw[very thick,SurfaceGreen] (0,-2.2)--(0,2.2);
    \draw[white,line width=5pt] (-2.5,0)--(2.5,0);
    \draw[very thick,DictionaryBlue,dashed] (-2.5,0)--(2.5,0);
    \node[font=\bfseries] at (0,3.0)
      {\(D(y)\subseteq\AAA^2\)};
    \node[align=center] at (1.35,1.15)
      {\(y\neq0\)\\
       \(\AAA^1\times\Gm\)};
  \end{scope}
\end{tikzpicture}
\caption{The face \(\tau=\Cone{e_1}\) gives the affine open subset
\(U_\tau=D(y)\subseteq U_\sigma=\AAA^2\).}
\label{fig:face-gives-affine-open}
\end{figure}

Faces are therefore the combinatorial instructions for passing to
torus-invariant affine open subsets. This is exactly the mechanism
needed to glue affine toric varieties: when two cones meet in a common
face, the corresponding affine varieties contain a common affine
open subset.

\begin{threedictionaries}
The main properties of \(U_\sigma\) can now be read in three
languages.

\medskip

\begin{center}
\small
\begin{tabularx}{\linewidth}{@{}
  >{\raggedright\arraybackslash}X
  >{\raggedright\arraybackslash}X
  >{\raggedright\arraybackslash}X@{}}
\toprule
\textbf{Algebra}
&
\textbf{Geometry}
&
\textbf{Combinatorics}
\\
\midrule
semigroup maps
\(\sigma^\vee\cap M\to\Bbbk\)
&
points of \(U_\sigma\)
&
multiplicative assignments to lattice points
\\[1mm]
units \(c\character{m}\)
&
torus factors
&
\(m\in\sigma^\perp\cap M\)
\\[1mm]
the maximal monomial ideal
&
the unique fixed point
&
\(\sigma\) is full-dimensional
\\[1mm]
a polynomial--Laurent coordinate ring
&
smoothness
&
primitive ray generators extend to a basis
\\[1mm]
pullback of characters
&
toric morphisms
&
\(F_{\RR}(\sigma_1)\subseteq\sigma_2\)
\\[1mm]
localization at \(\character{m}\)
&
a torus-invariant affine open subset
&
a face \(\tau=\sigma\cap m^\perp\)
\\
\bottomrule
\end{tabularx}
\end{center}
\end{threedictionaries}

\begin{guidingquestion}
How can we assemble the constructions and properties of this chapter
into one classification dictionary, and what do the standard examples
look like in all three languages?
\end{guidingquestion}


\section{The classification dictionary}
\label{sec:classification-affine-toric-varieties}
\index{classification!normal affine toric varieties}
\index{affine toric variety!classification}

We can now assemble the chapter into one statement. Fix an algebraic
torus \(T\), with lattices
\[
M=X^*(T),
\qquad
N=X_*(T).
\]

\begin{theorem}[Classification of normal affine toric varieties]
\label{thm:classification-normal-affine-toric-varieties}
The following data determine one another:
\begin{enumerate}[label=(\alph*),wide=0pt]
\item a normal affine toric variety \(X\) with dense torus \(T\);
\item a saturated affine semigroup \(S\subseteq M\) satisfying
\[
\ZZ S=M;
\]
\item a full-dimensional rational polyhedral cone
\[
C\subseteq M_{\RR};
\]
\item a strongly convex rational polyhedral cone
\[
\sigma\subseteq N_{\RR}.
\]
\end{enumerate}
The correspondences are
\[
S=C\cap M,
\qquad
C=\sigma^\vee,
\qquad
\sigma=C^\vee,
\]
and
\[
X
\simeq
\Spec\Bbbk[S]
=
\Spec\Bbbk[\sigma^\vee\cap M].
\]
\end{theorem}

\noindent
\textit{Reference.}
This combines Cox--Little--Schenck, Theorems~1.2.18 and~1.3.5.

\begin{proof}[Conceptual proof sketch]
Starting with \(\sigma\), Gordan's lemma makes
\[
S=\sigma^\vee\cap M
\]
finitely generated. It is saturated because it contains every lattice
point of its cone, so \(\Bbbk[S]\) is normal. Strong convexity ensures
that \(S\) generates \(M\), and therefore
\[
U_\sigma=\Spec\Bbbk[S]
\]
has dense torus \(T\).

Conversely, the regular characters on a normal affine toric variety
form a saturated affine semigroup \(S\subseteq M\). Hence
\[
S=\Cone{S}\cap M.
\]
Taking
\[
C=\Cone{S},
\qquad
\sigma=C^\vee
\]
recovers the variety. Biduality shows that no information is lost.
\end{proof}

\begin{warningbox}
Without normality, the cone does not determine the variety: it forgets
the holes of the semigroup. It determines only the normalization
\[
\Spec\Bbbk[\Cone{S}\cap M_S].
\]
\end{warningbox}

\clearpage
\section*{What this chapter has accomplished}
\addcontentsline{toc}{section}{What this chapter has accomplished}
\markboth{What this chapter has accomplished}{What this chapter has accomplished}

Starting with the affine plane, we followed the chain
\[
\text{extending characters}
\longrightarrow
S\subseteq M
\longrightarrow
\Bbbk[S]
\longrightarrow
\Spec\Bbbk[S].
\]
For a general affine toric variety, the semigroup \(S\) is the
discrete data that remembers every monomial allowed in the coordinate
ring.

Normality removes the holes:
\[
S\text{ saturated}
\quad\Longleftrightarrow\quad
S=\Cone{S}\cap M_S
\quad\Longleftrightarrow\quad
\Bbbk[S]\text{ normal}.
\]
At that point the cone contains all the information. Passing to the
dual lattice gives the standard formula
\[
\boxed{
U_\sigma
=
\Spec\Bbbk[\sigma^\vee\cap M].
}
\]

The formula has two complementary readings:
\[
\boxed{
\begin{aligned}
m\in\sigma^\vee\cap M
&\Longleftrightarrow
\character{m}\text{ extends regularly to }U_\sigma,\\
u\in\sigma\cap N
&\Longleftrightarrow
\oneps{u}\text{ extends across }t=0.
\end{aligned}
}
\]
Equivalently, the dual cone records regular functions, while the
original cone records one-parameter motions whose limits exist.

The cone also reads the basic geometry of \(U_\sigma\): orthogonal
directions give torus factors, full dimension gives a fixed point,
lattice bases give smoothness, linear maps give toric morphisms, and
faces give affine open subsets.

\begin{figure}[ht]
\centering
\begin{tikzpicture}[
  mainbox/.style={
    rounded corners,
    draw=black!55,
    text width=3.25cm,
    minimum height=1.55cm,
    align=center,
    inner sep=5pt
  }
]
  \node[mainbox,fill=DictionaryBlue!8] (N) at (-4.4,0)
    {\(\displaystyle
      \sigma\subseteq N_{\RR}\)\\[1mm]
     one-parameter directions\\
     whose limits exist};

  \node[mainbox,fill=PastelYellow] (M) at (0,0)
    {\(\displaystyle
      \sigma^\vee\cap M\)\\[1mm]
     characters\\
     extending regularly};

  \node[mainbox,fill=SurfaceGreen!9] (X) at (4.4,0)
    {\(\displaystyle
      U_\sigma=
      \Spec\Bbbk[\sigma^\vee\cap M]\)\\[1mm]
     normal affine\\
     toric variety};

  \draw[very thick,<->,PredictionTeal]
    (N.east)--(M.west);
  \draw[very thick,-{Latex[length=2.4mm]},DictionaryBlue]
    (M.east)--(X.west);

  \node[above=0.25cm of N] {motions};
  \node[above=0.25cm of M] {functions};
  \node[above=0.25cm of X] {geometry};
\end{tikzpicture}
\caption{The classification dictionary for normal affine toric
varieties.}
\label{fig:affine-toric-classification-dictionary}
\end{figure}

\begin{prediction}
A single cone produces one affine toric variety. Suppose several cones
meet along common faces. Since faces give common affine open subsets,
how might the varieties \(U_\sigma\) be glued to produce nonaffine
toric varieties?
\end{prediction}

\begin{guidingquestion}
How can compatible cones be glued along their common faces to produce
general normal toric varieties?
\end{guidingquestion}

\clearpage

\begin{exercises}[Comprehension exercises]
\markboth{Comprehension exercises}{Comprehension exercises}

\item Let
\[
S=\NN(2,0)+\NN(0,1)+\NN(1,1)\subseteq\ZZ^2.
\]
List all points of \(S\) of total degree at most \(4\), where total
degree means the sum of the coordinates. Compute \(\ZZ S\) and
\(\Cone{S}\).

\item For each of the following affine semigroups, write down the
semigroup algebra \(\Bbbk[S]\) in familiar coordinates and identify
\(\Spec\Bbbk[S]\):
\[
\NN^2,\qquad
\ZZ_{\geq0}\times\ZZ,\qquad
\ZZ^2.
\]

\item Let
\[
\mathcal A=\{(1,0),(1,1),(0,1)\}\subseteq\ZZ^2.
\]
Write the monomial map
\[
\Phi_{\mathcal A}:\Gm^2\longrightarrow\AAA^3.
\]
Find a binomial equation for the closure of its image and identify the
resulting affine toric variety.

\item Let
\[
S=\NN\{2,3\}\subseteq\ZZ.
\]
Show that \(S\) is not saturated. Compute \(\Cone{S}\cap\ZZ\), the
normalization of \(\Bbbk[S]\), and the normalization map of
\(\Spec\Bbbk[S]\).

\item Let
\[
S=\NN(2,0)+\NN(0,1)+\NN(1,1)\subseteq\ZZ^2.
\]
Show that
\[
S\subsetneq\Cone{S}\cap\ZZ^2.
\]
Find a point witnessing the failure of saturation, and compute the
normalization of \(\Bbbk[S]\).

\item Compute the dual cones of the following cones in \(\RR^2\):
\[
\Cone{(1,0)},\qquad
\Cone{(1,0),(0,1)},\qquad
\Cone{(1,0),(1,2)}.
\]
For each cone \(\sigma\), describe the semigroup
\[
\sigma^\vee\cap M.
\]

\item Let
\[
\sigma=\Cone{(1,0),(0,1)}\subseteq N_{\RR}\simeq\RR^2.
\]
For each lattice point
\[
u=(a,b)\in N,
\]
determine whether the limit
\[
\lim_{t\to0}\oneps{u}(t)
\]
exists in \(U_\sigma=\AAA^2\). When it exists, determine the limit
point.

\item Let
\[
\sigma=\Cone{(1,0),(1,2)}.
\]
Using the Hilbert basis of \(\sigma^\vee\cap M\), write the coordinate
ring of \(U_\sigma\) as a quotient of a polynomial ring. Show that
\(U_\sigma\) is normal but singular.

\item For each of the following cones in \(\ZZ^2\), determine whether
the associated affine toric variety is smooth:
\[
\Cone{(1,0),(0,1)},\qquad
\Cone{(1,0),(1,3)},\qquad
\Cone{(1,1),(-1,0)}.
\]
In the smooth cases, identify the variety as
\[
\AAA^2,\qquad
\AAA^1\times\Gm,\qquad
\text{or }\Gm^2.
\]

\item Let
\[
F:\ZZ^2\longrightarrow\ZZ,
\qquad
(a,b)\longmapsto a+b.
\]
Show that
\[
F_{\RR}\bigl(\Cone{(1,0),(0,1)}\bigr)\subseteq\RR_{\geq0},
\]
and use the cone criterion to produce a toric morphism
\[
\AAA^2\longrightarrow\AAA^1.
\]
Write the induced map on coordinate rings and the corresponding map of
varieties.

\item Let
\[
\sigma=\Cone{e_1,e_2}\subseteq\RR^2,
\qquad
\tau=\Cone{e_1}.
\]
Choose a lattice point \(m\in\sigma^\vee\cap M\) such that
\[
\tau=\sigma\cap m^\perp.
\]
Then verify directly that
\[
U_\tau
\simeq
D(\character{m})
\subseteq
U_\sigma.
\]

\end{exercises}

\clearpage
\section*{Challenge exercise}
\addcontentsline{toc}{section}{Challenge exercise}
\markboth{Challenge exercise}{Challenge exercise}

\noindent
\textbf{Cyclic quotient orbifolds and their coarse toric surfaces.}

\medskip

Assume that the characteristic of \(\Bbbk\) does not divide \(n\).
Choose an integer \(q\) satisfying
\[
1\leq q<n,
\qquad
\gcd(q,n)=1,
\]
and let \(\mu_n\) act on \(\AAA^2\) by
\[
\zeta\cdot(x,y)
=
(\zeta x,\zeta^q y).
\]

We write
\[
\left[\AAA^2/\mu_n\right]
\]
for the resulting smooth toric orbifold, and
\[
X
=
\AAA^2/\mu_n
=
\Spec\Bbbk[x,y]^{\mu_n}
\]
for its coarse moduli space.

No general theory of stacks is needed for this exercise. Think of
\(\left[\AAA^2/\mu_n\right]\) as the affine plane together with its
finite symmetry, and of \(X\) as the affine variety whose regular
functions are precisely the invariant polynomials.

The goal is to see that the orbifold and its coarse space use the
same real cone but different lattices. The ordinary coordinates
\(x,y\) will become \emph{orbifold pseudo-coordinates}: they are
honest coordinates upstairs, but usually not regular functions on the
coarse space.

Let
\[
\widetilde T=\Gm^2,
\qquad
\widetilde M=X^*(\widetilde T)\simeq\ZZ^2,
\qquad
\widetilde N=X_*(\widetilde T)\simeq\ZZ^2,
\]
and let
\[
\sigma=\Cone{e_1,e_2}\subseteq\widetilde N_{\RR}.
\]
Thus
\[
\AAA^2=U_{\sigma,\widetilde N}.
\]

\begin{enumerate}[label=\textbf{Hint \arabic*.},wide=0pt]

\item Show that the inclusion
\[
\mu_n\hookrightarrow\widetilde T,
\qquad
\zeta\longmapsto(\zeta,\zeta^q),
\]
recovers the given action on \(\AAA^2\). Explain why the coordinate
axes and the dense torus are preserved.

\item A character
\[
\character{(a,b)}=x^ay^b
\]
of \(\widetilde T\) restricts to the character
\[
\zeta\longmapsto\zeta^{a+qb}
\]
of \(\mu_n\). Show that the characters descending to the quotient
torus
\[
T=\widetilde T/\mu_n
\]
form the sublattice
\[
M
=
\left\{
(a,b)\in\widetilde M:
a+qb\equiv0\pmod n
\right\}.
\]
Show that
\[
[\widetilde M:M]=n.
\]

\item Prove that
\[
\Bbbk[x,y]^{\mu_n}
=
\bigoplus_{\substack{a,b\geq0\\
a+qb\equiv0\!\!\!\pmod n}}
\Bbbk\,x^ay^b.
\]
Deduce that
\[
\Bbbk[x,y]^{\mu_n}
=
\Bbbk[\sigma^\vee\cap M].
\]

\item Let
\[
N=\Hom_{\ZZ}(M,\ZZ).
\]
Dualizing the finite-index inclusion
\[
M\subseteq\widetilde M
\]
produces a finite-index inclusion
\[
\widetilde N\subseteq N.
\]
Show that
\[
N
=
\widetilde N
+
\ZZ\left(\frac1n(1,q)\right)
\subseteq\widetilde N_{\QQ}.
\]
In particular,
\[
[N:\widetilde N]=n.
\]

\item Regard the same real cone
\[
\sigma=\Cone{e_1,e_2}
\]
inside
\[
N_{\RR}=\widetilde N_{\RR}.
\]
Show that the coarse moduli space is the affine toric variety
\[
X=U_{\sigma,N}
=
\Spec\Bbbk[\sigma^\vee\cap M].
\]
Explain why the inclusion of lattices
\[
\widetilde N\hookrightarrow N
\]
induces the quotient morphism
\[
\AAA^2=U_{\sigma,\widetilde N}
\longrightarrow
U_{\sigma,N}=X.
\]

\item The coordinate functions \(x\) and \(y\) correspond to the
characters
\[
(1,0),(0,1)\in\widetilde M.
\]
Determine when either of these characters belongs to \(M\). Explain
why \(x\) and \(y\) are generally not functions on the coarse space,
even though invariant monomials in \(x\) and \(y\) are.

Interpret \(x,y\) as orbifold pseudo-coordinates.

\item Show that \(e_1,e_2\) are primitive lattice vectors in \(N\),
but that they span the proper sublattice
\[
\widetilde N\subseteq N
\]
of index \(n\). Conclude from the smoothness criterion that, for
\(n>1\),
\[
U_{\sigma,N}
\]
is singular at its torus-fixed point.

Contrast this with the smooth orbifold
\[
[\AAA^2/\mu_n].
\]
The same real cone is smooth in the lattice \(\widetilde N\), but not
in the larger lattice \(N\).

\item Specialize to
\[
q=n-1,
\]
so that
\[
\zeta\cdot(x,y)=(\zeta x,\zeta^{-1}y).
\]
Show that every invariant monomial is generated by
\[
u=x^n,
\qquad
v=y^n,
\qquad
w=xy.
\]
Find the relation among \(u,v,w\), and conclude that
\[
X
\simeq
V(uv-w^n)\subseteq\AAA^3.
\]

This is the \(A_{n-1}\) surface singularity.

\item Summarize the two toric descriptions:
\[
\AAA^2=U_{\sigma,\widetilde N},
\qquad
\AAA^2/\mu_n=U_{\sigma,N}.
\]
Explain in your own words how changing the lattice while keeping the
same real cone turns genuine coordinates upstairs into orbifold
pseudo-coordinates and creates a singular coarse moduli space.
\end{enumerate}


\clearpage

\section*{Notes and references}
\addcontentsline{toc}{section}{Notes and references}
\markboth{Notes and references}{Notes and references}

This chapter follows the first half of Cox--Little--Schenck,
Chapters~1 and~3, with the emphasis shifted toward the three
dictionaries:
\[
\text{algebra},\qquad
\text{geometry},\qquad
\text{combinatorics}.
\]
For the general theory of affine semigroups, Hilbert bases, and
normality, standard complementary references include Bruns--Gubeladze,
\emph{Polytopes, Rings, and K-Theory}, and Miller--Sturmfels,
\emph{Combinatorial Commutative Algebra}.

From the toric viewpoint, the central result is that normal affine
toric varieties are classified by strongly convex rational polyhedral
cones. The nonnormal examples in this chapter are included to show
exactly where cones stop being enough and semigroups must take over.

The next chapter globalizes this picture. Faces already provide common
affine open subsets, so compatible cones can be glued. This leads from
single cones to fans, and from affine toric varieties to general
normal toric varieties.
